Мы убеждены, что системы типов должны уметь
переинтерпретировать собственные термины в произвольной
внутренней модели: вооружившись единственным
инструментом --- реинтерпретацией --- можно выразить
широкий круг возможностей, таких как полиморфизм
вселенных, внешняя параметричность, возможность вести
синтетическую математику непосредственно внутри системы
и многое другое. В этой связи мы исследуем пределы
самореинтерпретации, представляем две системы типов,
способные к ней (HM-R и MLTT-R), доказываем их базовые
свойства и приводим несколько примеров полезных
внутренних моделей.

\section{Введение}

Рассмотрим следующую хорошо известную функцию,
применяющую заданное преобразование к элементам списка,
написанную на Agda:

\begin{minted}{agda}
map : ∀ {A B : Set} → (A → B) → List A → List B
map f [] = []
map f (x ∷ xs) = f x ∷ map f xs
\end{minted}

Наслаждаясь и композиционностью этой функции, и
простотой её типа, мы хотели бы максимально
переиспользовать её определение. В некоторой степени
это было бы, безусловно, возможно: ведь \texttt{map}
полиморфна относительно исходного и целевого типов, так
что мы можем подставлять туда любые желаемые типы.
Однако, если бы нам попалась волшебная лампа с джинном
теории типов, мы попросили бы его исполнить три
желания:

\begin{enumerate}
    \item Переиспользовать \texttt{map} на более
        высоких уровнях вселенных: например, имея
        \textit{оператор типов}, применить map к
        \textit{списку типов}:
        \[
            (\mathtt{Set}\to\mathtt{Set})\to
            \mathtt{List}\,\mathtt{Set}\to
            \mathtt{List}\,\mathtt{Set}
        \]
    \item Автоматически получить доказательство
        свободной теоремы для \texttt{map}, то есть
        некоторый терм типа
        \begin{multline*}
            \forall A\,A'\,(Q:A\to A'\to\mathtt{Prop}),
            \forall B\,B'\,(R:B\to B'\to\mathtt{Prop}),
            \\
            \forall f\,f',\left(
            \forall x\,x',Q\,x\,x'\to
            R(f x)\left(f' x'\right)
            \right)\to\\
            \forall
            (\mathtt{xs}:\mathtt{List}\,A)
            (\mathtt{xs}':\mathtt{List}\,A'),
            Q^*\,\mathtt{xs}\,\mathtt{xs}'\to
            R^*(\mathtt{map}\,f\,\mathtt{xs})
            (\mathtt{map}\,f'\,\mathtt{xs}');
        \end{multline*}
    \item По данным группоидам $A$ и $B$ и функтору
        $f: A \to B$ построить функтор
        $\mathtt{map}\;f$ между группоидами
        $\mathtt{List}\,A$ и $\mathtt{List}\,B$.
\end{enumerate}

Первое желание, разумеется, исполняется: достаточно
доработать определение, сделав его полиморфным
относительно уровней. Хотя теорема из второго желания
уже верна в системах с параметричностью, получить её
доказательство \textit{внутри} системы без участия
человека можно лишь в тех системах типов, где
параметричность интернализирована
\cite{Bernardy:Parametricity2010,
Altenkirch:Parametricity2024}. Третье желание пока
приходится исполнять вручную (или с привлечением
метапрограммирования \cite{Nawrocki:Lean2026}), хотя
для этого требуется лишь пара индукций по списку:
несколько для определения экземпляра группоида для
\texttt{List} и ещё несколько для доказательства
функториальности \texttt{map}.

Мы предлагаем иной подход. Прежде чем дорабатывать
определение \texttt{map}, вспомним всю
метатеоретическую мощь, которой наша система типов уже
обладает. Первой на ум приходит упомянутая выше
параметричность полиморфных определений; обратим
внимание, что она получается путём задания
\textit{реляционной интерпретации} суждений выводимости
\cite{reynolds1983types}. Действительно, некоторые из
наиболее интересных метатеоретических результатов
выражаются через построение определённой
\textit{модели} системы типов: синтетические результаты
в теории гомотопий опираются на
$(\infty, 1)$-группоидную модель (гомотопической)
теории типов
\cite{HofmannStreicher-Thegroupoidmodelref,
altenkirch:groupoids}; применение Макбрайдом операции
смещения для полиморфизма вселенных
\cite{McBride:Crude2011,McBride:Crude2012} можно
рассматривать как построение модели системы типов в
более высоких вселенных её же самой, а затем
интерпретацию смещённых конструкций в ней. Тогда все
перечисленные выше желания могли бы быть исполнены
путём построения подходящей модели системы типов
\textit{внутри} неё же (отсюда наш термин
\textit{внутренние модели}) и последующего
инстанциирования \texttt{map} в ней. Ниже мы будем
называть системы типов с возможностью
самореинтерпретации \textit{рефлексивными}.

\subsection{Вклад}

\begin{itemize}
    \item Наша главная цель --- построить реализуемую
        рефлексивную систему типов с зависимыми типами.
        Основная трудность при этом возникает из-за
        результата о невозможности, обсуждаемого в
        разделе \ref{no-go-section}, который вынуждает
        нас ограничить мощность реинтерпретации, чтобы
        сохранить алгоритмичность проверки типов.
    \item Ради удобства изложения мы сначала предлагаем
        рефлексивную систему типов на основе системы
        типов Хиндли\,---\,Милнера, названную HM-R
        (раздел \ref{HM-R}). Мы доказываем, что HM-R
        корректна, консервативна относительно просто
        типизированного лямбда-исчисления, и что
        глобальный вывод типов алгоритмичен (раздел
        \ref{HM-R-properties}).
    \item Опираясь на наработки, полученные при
        разработке HM-R, мы наконец готовы погрузиться
        в рефлексивную версию теории типов
        Мартина-Лёфа, именуемую MLTT-R (раздел
        \ref{MLTT-R}). Мы также доказываем корректность
        и консервативность MLTT-R (раздел
        \ref{MLTT-R-properties}).
    \item Как и планировалось, благодаря
        реинтерпретации MLTT-R объединяет и
        ограниченную форму полиморфизма вселенных, и
        \textit{внутреннюю внешнюю} параметричность. И,
        разумеется, MLTT-R можно использовать как
        систему для синтетической математики с
        внутренним переиспользованием синтетических
        результатов (раздел \ref{MLTT-R-applications}).
\end{itemize}

\section
{Теорема о невозможности внутренней интерпретации}
\label{no-go-section}

Итак, мы хотим построить рефлексивную систему типов, то
есть систему, в которой пользователь может как
определить модель системы внутри неё, так и применить
её к корректно типизированным термам, получив их
внутренние интерпретации. Разумеется, нам хотелось бы,
чтобы определяемые пользователем модели сохраняли
дефинициальные равенства системы; в системе с
зависимыми типами это естественно достигается с помощью
типа тождества, позволяющего явно указать
соответствующие равенства внутри модели. Было бы также
хорошо, если бы перед реинтерпретацией терма система
сначала приводила его к нормальной форме, чтобы
избавиться от артефактов построения в результате
интерпретации. Наконец, в системе с зависимыми типами
естественно ожидать редукцию открытых термов, так что
мы хотим, чтобы наша интерпретация вычисляла значения
на открытых термах.

Доказываемая ниже теорема утверждает, что в одной
системе невозможно иметь все три свойства одновременно.

Сначала пара определений. Мы будем работать в рамках
классической типотеоретической схемы, которая вводит
контексты $\Gamma$, термы $t, t'$ и типы $T$, связанные
отношением типизации $\Gamma\vdash t:T$; вычисление
задаётся одношаговым отношением редукции $t\leadsto t'$
на термах.

\begin{definition}
    [Наивная вычислительная реинтерпретация]
    Оператор типов $F[\_]$ называется наивной
    вычислительной реинтерпретацией типов, если он
    сохраняет вычисления на типах (то есть, если на
    типах существует отношение редукции $(\leadsto)$ и
    $T\leadsto T'$, то $F[T]\leadsto^* F[T']$) и
    вычисляется как примитивно рекурсивная функция от
    нормальных форм типов. Аналогично, терм с дыркой
    $\llbracket\_\rrbracket$ называется наивной
    вычислительной реинтерпретацией термов, если он
    сохраняет вычисления на термах, вычисляется как
    примитивно рекурсивная функция от нормальных форм
    термов и, если $\Gamma\vdash t:T$, то
    $\Gamma\vdash\llbracket t\rrbracket:F[T]$.
\end{definition}

\begin{theorem}
    [Теорема о невозможности наивной реинтерпретации]
    \label{no-go}
    В конфлюэнтной системе типов с типом функций и
    наивной вычислительной реинтерпретацией, если
    редукция коммутирует с подстановкой, то для
    произвольных корректно типизированных $s$ и $t$
    существует $u$ такой, что:
    \[
        \mathtt{app}\llbracket\lambda x.t\rrbracket
        \llbracket s \rrbracket \leadsto^* u\;
        \text{\reflectbox{$\leadsto^*$}}\;
        \llbracket t[x\coloneqq s] \rrbracket
    \]
    где $\mathtt{app}$ — действие
    $\llbracket\_\rrbracket$ на аппликациях.
\end{theorem}

Следовательно,
$\mathtt{app}\llbracket\lambda x.t\rrbracket
\llbracket s \rrbracket$ и
$\llbracket t[x\coloneqq s]\rrbracket$ дефинициально
эквивалентны, что вынуждает бета-редукции в модели
выполняться дефинициально, а не пропозиционально, что
прямо противоречит нашему желаемому результату.
Аналогичные теоремы можно доказать для любой другой
редукции, присутствующей в системе; иными словами,
\textbf{пропозициональное равенство недостаточно
сильно}. В качестве доказательства приведём конкретный
пример:

\begin{proof}
    Даны корректно типизированные $\Gamma\vdash s:S$ и
    $\Gamma,x:S\vdash t:T$. Рассмотрим следующий терм:
    \[
        \Gamma,x:S\to T,y:S\vdash
        \llbracket x\,y\rrbracket:F[T]
    \]
    Терм под интерпретацией находится в нормальной
    форме, следовательно, весь терм вычисляется в
    $\mathtt{app}\llbracket x\rrbracket
    \llbracket y\rrbracket$. Подставляя $\lambda x.t$
    вместо $x$ и $s$ вместо $y$, получаем
    $\mathtt{app}\llbracket\lambda x.t\rrbracket
    \llbracket s \rrbracket$.

    С другой стороны, если сначала выполнить
    подстановку, мы получили бы
    $\Gamma\vdash\llbracket(\lambda x.t)s
    \rrbracket:F[T]$. Поскольку наивная реинтерпретация
    сохраняет редукцию, это редуцируется в
    $\llbracket t[x\coloneqq s] \rrbracket$.

    Поскольку редукция коммутирует с подстановкой, а
    система конфлюэнтна, должно существовать такое $u$,
    в которое редуцируются оба терма.
\end{proof}

Доказательство выше эксплуатирует открытость исходного
терма; если мы ограничим реинтерпретацию вычислением
только на замкнутых термах, мы сможем получить
корректную систему без требования дефинициальных
равенств в моделях, как это демонстрирует наша система
без зависимых типов HM-R, описанная в разделе
\ref{HM-R}. Однако для системы с зависимыми типами наша
теорема о невозможности оставляет нам три варианта:

\begin{enumerate}
    \item Ограничить модели теми, в которых равенства
        выполняются дефинициально, а не
        пропозиционально.
    \item Не выполнять реинтерпретацию, пока терм не
        является замкнутым. Для рассуждений над
        открытыми терминами уравнения, описывающие
        поведение интерпретации, выставляются как
        пропозициональные аксиомы.
    \item Не редуцировать открытые термы под
        реинтерпретацией вовсе, выполняя интерпретацию 
        как есть. Для упрощения результата
        реинтерпретации можно переписывать его с
        помощью равенств, присутствующих в модели.
\end{enumerate}

Если рассматривать существующие работы о дисплейсментах
\cite{McBride:Crude2011,McBride:Crude2012}, внутренней
параметричности \cite{Altenkirch:Parametricity2024} и
отображённой теории типов
\cite{Kolomatskaia_Shulman_2025} через призму
реинтерпретации внутри фиксированной модели (или
семейства моделей), то видно, что соответствующие
модели действительно выполняют равенства дефинициально;
следовательно, по-видимому, необходимо принять первый
подход. К тому же он наиболее удовлетворителен с
вычислительной точки зрения. Однако, если Вам нужна
корректная система, в которой модели могут определяться
пользователем, необходимо ограничивать модели через
дисциплину типизации. Используя известные на данный
момент методики
\cite{bezem2021notegeneralizedalgebraictheories},
такое ограничение моделей потребовало бы определения
типа моделей как типа алгебр
индуктивно-индуктивно-рекурсивного типа (или даже
индуктивно-индуктивно-рекурсивного симплициального
типа, для унивалентных теорий типов), что значительно
выходит за рамки данной статьи.

Из двух оставшихся вариантов последний лучше подходит
для практических систем с зависимыми типами, поскольку
он надлежащим образом вычисляет значения на открытых
термах. К тому же он не ограничивает класс возможных
моделей, при этом предоставляя преимущества первого
подхода: если модель такова, что уравнения выполняются
дефинициально, то переписывание результата
реинтерпретации для его редукции не требуется.
Единственный недостаток состоит в том, что, возможно,
редукции \textit{после} интерпретации обходятся дороже,
чем те, которые мы могли бы выполнить \textit{до} неё.

\section
{Почти просто типизированный почти рефлексивный язык:
HM-R}\label{HM-R}

Мы начинаем с более простого языка, чтобы проработать
все детали и разобраться, как правильно
реинтерпретировать термы, прежде чем переходить к
системе с зависимыми типами. Простейший язык, для
которого нам известно множество интересных моделей, ---
это просто типизированное лямбда-исчисление с
произведениями, которое мы в течение всей статьи будем
называть ``STLC''; в этом разделе мы представим язык,
способный интерпретировать термы STLC как конструкции в
моделях STLC, которые могут быть определены в нашем
языке.

Итак, что же представляют собой модели STLC? Помимо
обычной теоретико-множественной семантики, где $S\to T$
интерпретируется как функция из $S$ в $T$, а
$S\times T$ --- как декартово произведение, существует
теоретико-доказательственная семантика, которая
связывает термы типа $T$ с интуиционистскими
доказательствами пропозиции $\llbracket T\rrbracket$,
где $(\to)$ теперь понимается как импликация, а
$(\times)$ --- как конъюнкция. Это, конечно,
классическое соответствие Карри--Говарда \cite{ho80};
обе семантики --- теоретико-множественная и
доказательственная --- обобщаются полным описанием
моделей STLC, данным Йоханом Ламбеком
\cite{Lambek1980}. В соответствии с этим
\textit{соответствием Карри--Говарда--Ламбека}, модели
STLC суть именно декартово замкнутые категории. Им
можно дать множество различных определений; мы будем
использовать определение \ref{ccc}.

\begin{definition}[Декартово замкнутая категория]
    \label{ccc}
    Декартово замкнутая категория --- это следующие
    данные:
    \begin{enumerate}
        \item Множество \textit{объектов} $\mathrm{Ob}$
            и семейство \textit{морфизмов}
            $\mathrm{Hom}(S, T)$, индексированное
            $S, T\in\mathrm{Ob}$, где
            $f\in\mathrm{Hom}(S,T)$ обозначается как
            $f:S\to T$. Кроме того, $S$ называется
            \textit{исходным объектом} $f$, а $T$ ---
            \textit{целевым объектом} $f$.
        \item Для каждого $A\in\mathrm{Ob}$ существует
            \textit{тождественный морфизм}
            $\mathrm{id}_A: A\to A$; для каждого
            $f:A\to B$ и $g:B\to C$ существует их
            \textit{композиция} $g\circ f:A\to C$.
            Если объект-цель тождественного морфизма
            ясен из контекста, нижний индекс
            опускается. Кроме того, для произвольных
            морфизмов $f, g, h$ с согласованными
            исходными и целевыми объектами выполняются
            следующие свойства:
            \begin{align*}
                \mathrm{id}\circ f = f =
                f\circ\mathrm{id} & &
                (f \circ g) \circ h =
                f \circ (g \circ h)
            \end{align*}
        \item Существует \textit{терминальный объект}
            $\top\in\mathrm{Ob}$ такой, что для каждого
            $A\in\mathrm{Ob}$ существует
            \textbf{единственный} морфизм
            $!_A:A\to\top$.
        \item Для каждого $A, B\in\mathrm{Ob}$
            существует их \textit{произведение}
            $A\times B$ вместе с \textit{проекциями}
            $\pi_1:A\times B\to A$ и
            $\pi_2:A\times B\to B$; для каждого объекта
            $C\in\mathrm{Ob}$ и пары морфизмов
            $f:C\to A$ и $g:C\to B$ существует
            \textbf{единственный} \textit{морфизм пары}
            $[f;g]:C\to A\times B$ такой, что
            выполняются следующие свойства:
            \begin{align*}
                \pi_1\circ[f;g]=f & & \pi_2\circ[f;g]=g
            \end{align*}
        \item Для каждого $A, B\in\mathrm{Ob}$
            существует их \textit{экспоненциальный
            объект} $B^A$ вместе с \textit{морфизмом
            применения} $\mathrm{ev}:B^A\times A\to B$;
            для каждого объекта $C\in\mathrm{Ob}$ и
            морфизма $f:C\times A\to B$ существует
            \textbf{единственный} \textit{морфизм
            абстракции} $\lambda\,f:C\to B^A$ такой,
            что для произвольного $g:C\to A$
            выполняется следующее свойство:
            \[
                \mathrm{ev}\circ[\lambda\,f; g]
                = f\circ[\mathrm{id}_C; g]
            \]
    \end{enumerate}
\end{definition}

И тогда корректно типизированный терм
$\Gamma\vdash t:T$ интерпретируется как морфизм
$\llbracket t\rrbracket:
\llbracket\Gamma\rrbracket\to\llbracket T\rrbracket$,
где $\llbracket\Gamma\rrbracket$ --- вложенное влево
произведение интерпретаций типов из $\Gamma$:
\[
    \llbracket x:A,y:B,z:C\rrbracket=
    ((\top\times\llbracket A\rrbracket)
    \times\llbracket B\rrbracket)
    \times\llbracket C\rrbracket.
\]
Тип единицы теперь --- терминальный объект;
произведение типов теперь --- произведение объектов;
тип стрелки --- экспоненциальный объект. Уравнения в
(2) необходимы для выполнения структурных свойств
подстановки; уравнения в (4) соответствуют редукциям
пар; уравнение в (5) выполняет $\beta$-редукцию.
Свойства \textbf{уникальности} соответствуют
$\eta$-эквивалентности.

Например, в нашем языке
$\llbracket t\rrbracket\,\mathcal{M}$ --- это оператор,
реинтерпретирующий терм $t$ во внутренней модели
$\mathcal{M}$. Для следующего терма:
\[
\cdot\vdash\lambda f x y. (f x, y):
(A\to B)\to A\to C\to B\times C
\]
вычисление происходит следующим образом:
\begin{multline*}
    \llbracket\lambda f x y. (f x, y))\rrbracket\,
    \mathcal{M}=\\=
    \mathtt{lam}\,(\mathtt{lam}\,(\mathtt{lam}\,(
    \mathtt{pair}\,
    (\mathtt{eval}\,\circ\,\mathtt{pair}\,
    (\mathtt{proj}_2\circ\mathtt{proj}_1
    \circ\mathtt{proj}_1
    )\,
    (\mathtt{proj}_2\circ\mathtt{proj}_1)
    )\,
    \mathtt{proj}_2
    )))
\end{multline*}

где $\mathtt{lam}$, $\mathtt{pair}$ и $\mathtt{eval}$
--- морфизм абстракции, морфизм пары и морфизм
применения, взятые из $\mathcal{M}$ и определённые
таким же образом, как в определении \ref{ccc},
переведённые в типы нашей системы.

Теперь перейдём к самому проектированию нашего языка.
Как обсуждалось в разделе \ref{no-go-section}, для
независимого случая, с которым мы сейчас имеем дело, мы
будем интерпретировать только нормальные формы
замкнутых терминов. В частности, это означает, что нам
фактически не нужно, чтобы уравнения и
\textbf{уникальности} из определения \ref{ccc}
выполнялись в наших внутренних моделях! Такой подход не
нарушает ни одного из метатеоретических свойств,
рассмотренных в разделе \ref{HM-R-properties}. Кроме
того, для простоты примера мы будем рассматривать
только модели, в которых объекты --- это типы. Тогда
интернализация определения \ref{ccc} потребует
1)~операторов типов для определения типов морфизмов,
произведений и экспоненциальных объектов;
2)~полиморфных функций для операционализации
$\forall\exists$ конструкций в модели. На самом деле,
мы от операторов типов тоже избавимся, поскольку нам
достаточно лишь типов в контексте, которые затем
подставляются в определение модели для проверки
операций на уровне терминов. Таким образом, наш язык
представляет собой небольшое расширение классической
системы типов Хиндли\,---\,Милнера, отсюда и название
HM-R; грамматические правила собраны на рис.
\ref{hm-r-grammar}. Обратите внимание, что, несмотря на
наше описание, грамматика не определяет оператор
реинтерпретации $\llbracket t\rrbracket\,\mathcal{M}$
для произвольного термина $t$; вместо этого
реинтерпретации подлежат только переменные. В этой
системе реинтерпретация --- не что иное, как полная
функция на замкнутых значениях.

\begin{definition}[Значения в HM-R]
    Значения обозначаются метапеременными $u, v$ в
    BNF-грамматике ниже. Метапеременная $n$ обозначает
    нейтральные термы:
    \begin{align*}
        u,v&\Coloneqq
            n\;|\;()\;|\;(u,v)\;|\;\lambda x.v\\
        n&\Coloneqq x\;|\;\pi_i\,n\;|\;n\,v
    \end{align*}
\end{definition}

\begin{figure}
    \centering
    \begin{bnf}
        $S,T,C,P,E,TT$ : \textsf{MonoType} ::=
        | $\alpha$ : переменная типа
        | $\top$ : тип единицы
        | $S\times T$ : произведение типов
        | $S\to T$ : тип стрелки
        ;;
        $\sigma$ : \textsf{PolyType} ::=
        | $T$ : моноформный тип
        | $\forall\alpha.\sigma$ : универсальный тип
        ;;
        $s,t$ : \textsf{Term} ::=
        | $x$ : переменная
        | $()$ : единичное значение
        | $(s, t)$ : пара
        | $\pi_1\,t$ : первая проекция
        | $\pi_2\,t$ : вторая проекция
        | $\lambda x. t$ : абстракция
        | $t\,s$ : применение
        | \texttt{let}\,$x=s$\,\texttt{in}\,$t$
        : полиморфное связывание
        | $x^\mathcal{M}$ : реинтерпретация
        ;;
        $\mathcal{M}$ : \textsf{Model} ::=
        | $\langle T, t, \ldots \rangle$
        : определение модели (см. правила типизации)
    \end{bnf}
    \caption{Типы и термы HM-R}
    \label{hm-r-grammar}
\end{figure}

Таким образом, для определения реинтерпретации
достаточно рассмотреть термы STLC:

\begin{definition}[Реинтерпретация термов HM-R]
    \begin{align}{}
        \llbracket()\rrbracket^\mathcal{M}/\gamma &=
        \mathtt{cnst}
        \\
        \llbracket(s,t)\rrbracket^\mathcal{M}/\gamma &=
        \mathtt{pair}
        \,\llbracket s\rrbracket^\mathcal{M}/
        \gamma
        \,\llbracket t\rrbracket^\mathcal{M}/
        \gamma
        \\
        \llbracket\pi_i\,t\rrbracket^\mathcal{M}/\gamma
        &=\mathtt{comp}\,\mathtt{proj}_i\,\llbracket
        t\rrbracket^\mathcal{M}/\gamma
        \\
        \llbracket x\rrbracket^\mathcal{M}/\gamma &=
        \gamma\,x
        \\
        \llbracket\lambda x.t\rrbracket^\mathcal{M}/
        \gamma &= \mathtt{lam}\,\llbracket t
        \rrbracket^\mathcal{M}/\left[
        \vec{y}\coloneqq\mathtt{comp}\,
        \mathtt{proj}_1(\gamma\,\vec{y}),\;
        x\coloneqq\mathtt{proj}_2
        \right]
        \\
        \llbracket t\,s\rrbracket^\mathcal{M}/\gamma &=
        \mathtt{comp}\,\mathtt{eval}\,\left(
        \mathtt{pair}
        \,\llbracket t\rrbracket^\mathcal{M}/
        \gamma
        \,\llbracket s\rrbracket^\mathcal{M}/
        \gamma
        \right)
    \end{align}
    где \texttt{cnst},\,\texttt{pair} и другие операции
    модели взяты из $\mathcal{M}$, а $\gamma$ ---
    подстановка.
\end{definition}

Вот наш подход к обеспечению принципа, согласно
которому \textit{интерпретируются только замкнутые
термы}: как только терм $s$, связанный полиморфным
let-связыванием $\mathtt{let}\,x=s\,\mathtt{in}\,t$,
становится замкнутым и приводится к нормальной форме,
только тогда он подставляется в $t$, а всякий раз,
когда он подставляется в $x^\mathcal{M}$, он немедленно
переинтерпретируется. Эта операция обозначается как
$t\llbracket x\coloneqq s\rrbracket$ в наших правилах
вычисления, данных на рис.~\ref{hm-r-comp}. Как и
прежде, она ведёт себя как обычная подстановка во всех
случаях, за исключением переменных, находящихся под
переинтерпретацией:
\[
\left(x^\mathcal{M}\right)\llbracket x\coloneqq s
\rrbracket = s^\mathcal{M}/\varepsilon
\quad\text{($\varepsilon$~— пустая подстановка)}
\]
Легко заметить, что вычисление зависло бы, если бы
let-связанные термы содержали переменные, пришедшие из
неприменённых $\lambda$-абстракций. Наши правила
типизации, представленные на рис.~\ref{hm-r-typing},
защищают нас от этого, позволяя ссылаться только на
let-связанные термы внутри других let-связанных термов.
И вновь, помимо этого и правила типизации моделей,
здесь нет ничего отличного от стандартных декларативных
правил типизации системы типов Хиндли~---~Милнера
\cite{Clement:ML1986}.

\begin{figure}
    \centering
    \begin{mathpar}
        \inferrule{s\leadsto s'}{(s,t)\leadsto(s',t)}
        \and
        \inferrule{t\leadsto t'}{(s,t)\leadsto(s,t')}
        \and
        \inferrule{t\leadsto t'}{\pi_i t\leadsto\pi_i t'}
        \and
        \inferrule{}{\pi_i(t_1, t_2)\leadsto t_i}
        \\
        \inferrule{t\leadsto t'}
        {\lambda x.t\leadsto\lambda x.t'}
        \and
        \inferrule{s\leadsto s'}{s\,t\leadsto s'\,t}
        \and
        \inferrule{t\leadsto t'}{s\,t\leadsto s\,t'}
        \and
        \inferrule{}{(\lambda x.t)s\leadsto t[x\coloneqq s]}
        \\
        \inferrule{s\leadsto s'}{
            \mathtt{let}\,x=s\,\mathtt{in}\,t\leadsto
            \mathtt{let}\,x=s'\,\mathtt{in}\,t
        }
        \and
        \inferrule{t\leadsto t'}{
            \mathtt{let}\,x=s\,\mathtt{in}\,t\leadsto
            \mathtt{let}\,x=s\,\mathtt{in}\,t'
        }
        \and
        \inferrule{\mathrm{fv}(v)=\varnothing}{
            \mathtt{let}\,x=v\,\mathtt{in}\,t\leadsto
            t\llbracket x\coloneqq v\rrbracket
        }
    \end{mathpar}
    \caption{Правила вычислений в HM-R}
    \label{hm-r-comp}
\end{figure}

Чтобы проверять модели, мы немного злоупотребляем нашей
системой типов, чтобы не вводить операторов типов и не
усложнять её: типы морфизмов $C$, объектов-произведений
$P$ и экспоненциалов $E$ задаются как открытые
мономорфные типы, которые могут использовать две
типовые переменные, $0$ и $1$. Этот факт отражается в
синтаксисе модели с помощью нотации $\Lambda 01.C$,
$\Lambda 01.P$ и $\Lambda 01.E$ соответственно. $0$ и
$1$ затем подставляются при типизации операций модели
(нотация $T[S_1, S_2]$ является сокращением для
$T[0\coloneqq S_1,1\coloneqq S_2]$). Терминальный
объект~--- это просто тип с именем $TT$.

Чтобы проверять переинтерпретацию, нам нужно
определить, как модель действует на мономорфные типы.
Поскольку на них не выполняется вычисление,
переинтерпретация типов снова определяется как
рекурсивная функция от синтаксиса:

\begin{definition}[Переинтерпретация типов в~HM-R]
    Пусть $\mathcal{M}$~--- корректная модель, а
    $\sigma$~--- подстановка мономорфных типов. Тогда
    интерпретация типа в модели есть следующая функция:
    \begin{align}{}
        \llbracket T\rrbracket^\mathcal{M}_\sigma &=
        C[TT,T^\mathcal{M}_\sigma] \\
        \alpha^\mathcal{M}_\sigma &= \sigma\alpha \\
        \top^\mathcal{M}_\sigma &= TT \\
        (S\times T)^\mathcal{M}_\sigma &=
        P[S^\mathcal{M}_\sigma,T^\mathcal{M}_\sigma] \\
        (S\to T)^\mathcal{M}_\sigma &=
        E[S^\mathcal{M}_\sigma,T^\mathcal{M}_\sigma].
    \end{align}
\end{definition}

\begin{figure}
    \centering
    \begin{mathpar}
        \inferrule{}{\Gamma\vdash ():\top}
        \and
        \inferrule{\Gamma\vdash s:S\\\Gamma\vdash t:T}
        {\Gamma\vdash (s,t):S\times T}
        \and
        \inferrule{\Gamma\vdash t:S\times T}
        {\Gamma\vdash \pi_1\,t:S}
        \and
        \inferrule{\Gamma\vdash t:S\times T}
        {\Gamma\vdash \pi_2\,t:T}\\
        \and
        \inferrule{x:T\in\Gamma}{\Gamma\vdash x:T}
        \and
        \inferrule{\Gamma,x:S\vdash t:T}
        {\Gamma\vdash\lambda x.t:S\to T}
        \and
        \inferrule{\Gamma\vdash t:S\to T\\\Gamma\vdash s:S}
        {\Gamma\vdash t\,s:T}
        \and
        \inferrule{
            \forall\Gamma\vdash s:S\\
            \mathrm{fv}(S)=\vec{\alpha}\\
            \forall\Gamma,x:\forall\vec{\alpha}.S,
            \Delta\vdash t:T
        }{\forall\Gamma,\Delta\vdash
            \mathtt{let}\,x=s\,\mathtt{in}\,t:T}
        \and
        \inferrule{
            x:\forall\vec{\alpha}.T\in\forall\Gamma\\
            \mathcal{M}\,\mathrm{model}
        }{\forall\Gamma,\Delta\vdash x^\mathcal{M}
            :\llbracket T\rrbracket^\mathcal{M}_{
                \left[\vec{\alpha}\coloneqq
                    \vec{S}\right]}}
        \and
        \inferrule{
            \mathrm{fv}(C)\cup\mathrm{fv}(P)\cup
            \mathrm{fv}(E)\subseteq\{0,1\}\\
            \Gamma\vdash\mathtt{id}:C[\alpha,\alpha]\\
            \Gamma\vdash\mathtt{comp}:
            C[\beta,\gamma]\to C[\alpha,\beta]\to
            C[\alpha,\gamma]\\
            \mathrm{fv}(TT)=\varnothing\\
            \Gamma\vdash\mathtt{cnst}:C[\alpha,TT]\\
            \Gamma\vdash\mathtt{pair}:
            C[\gamma,\alpha]\to C[\gamma,\beta]\to
            C[\gamma,P[\alpha,\beta]]\\
            \Gamma\vdash\mathtt{proj}_1:
            C[P[\alpha,\beta],\alpha]\\
            \Gamma\vdash\mathtt{proj}_2:
            C[P[\alpha,\beta],\beta]\\
            \Gamma\vdash\mathtt{lam}:
            C[P[\gamma,\beta],\alpha]\to
            C[\gamma,E[\alpha,\beta]]\\
            \Gamma\vdash\mathtt{eval}:
            C[P[E[\alpha,\beta],\beta],\alpha]
        }{\Gamma\vdash\langle
            \Lambda 01.C,\mathtt{id},\mathtt{comp}, TT,
            \mathtt{cnst}, \Lambda 01.P,\mathtt{pair},
            \mathtt{proj}_1,\mathtt{proj}_2,
            \Lambda 01.E,\mathtt{lam},\mathtt{eval}
            \rangle\,\mathrm{model}}
    \end{mathpar}
    \caption{Декларативные правила типизации HM-R}
    \label{hm-r-typing}
\end{figure}

\subsection{Метатеоретические свойства HM-R}
\label{HM-R-properties}

Здесь мы рассматриваем лишь некоторые базовые свойства
системы, чтобы убедиться, что всё работает как
ожидается: вычисления конфлюэнтны; типы сохраняются при
любом вычислении; при необходимости все типы могут быть
выведены; всё терминирует, и с точностью до редукции
\texttt{let}-связываний это тот же самый STLC, с
которого мы начали. Доказательства в основном
используют стандартные приёмы, поскольку правила
переинтерпретации устроены так, что их можно напрямую
встроить в доказательства свойств STLC и системы типов
Хиндли~---~Милнера.

\begin{theorem}[Теорема Чёрча~---~Россера для~HM-R]
    Рефлексивное транзитивное замыкание $(\to)$
    обладает свойством Чёрча~---~Россера.
\end{theorem}
\begin{proof}
    Классическое доказательство через отношение
    параллельной редукции \cite{CRproof}. Параллельная
    редукция определяется следующим образом:
    \begin{mathpar}
    \inferrule{}{t\twoheadrightarrow t}
    \and
    \inferrule{
        s\twoheadrightarrow s'\\
        t\twoheadrightarrow t'
    }{(s,t)\twoheadrightarrow(s',t')}
    \and
    \inferrule{t\twoheadrightarrow t'}
    {\pi_i\,t\twoheadrightarrow\pi_i\,t'}
    \and
    \inferrule{
        t_1\twoheadrightarrow t_1'\\
        t_2\twoheadrightarrow t_2'
    }
    {\pi_i(t_1,t_2)\twoheadrightarrow t_i'}
    \\
    \inferrule{t\twoheadrightarrow t'}
    {\lambda x.t\twoheadrightarrow\lambda x.t'}
    \and
    \inferrule{
        s\twoheadrightarrow s'\\
        t\twoheadrightarrow t'
    }
    {s\,t\twoheadrightarrow s'\,t'}
    \and
    \inferrule{
        s\twoheadrightarrow s'\\
        t\twoheadrightarrow t'
    }
    {(\lambda x.s)\,t\twoheadrightarrow
        s'[x\coloneqq t']}
    \\
    \inferrule{
        s\twoheadrightarrow s'\\
        t\twoheadrightarrow t'
    }{
        \mathrm{let}\,x=s\,\mathrm{in}\,t
        \twoheadrightarrow
        \mathrm{let}\,x=s'\,\mathrm{in}\,t'
    }
    \and
    \inferrule{t\twoheadrightarrow t'}{
        \mathrm{let}\,x=v\,\mathrm{in}\,t
        \twoheadrightarrow
        t'\llbracket x\coloneqq v\rrbracket
    }
\end{mathpar}
\end{proof}

\begin{theorem}[Сохранение типов в~HM-R]
    Имеет место следующее:
    \begin{enumerate}
        \item Если $\mathcal{M}$~--- модель,
        $\gamma\,x:C[\Gamma^\mathcal{M},
        \Gamma(x)^\mathcal{M}]$ для всех $x\in\Gamma$
        и $\Gamma\vdash t:T$, то $\llbracket t
        \rrbracket^\mathcal{M}/\gamma:
        C[\Gamma^\mathcal{M},T^\mathcal{M}]$.
        \item Если $\Gamma\vdash s:S$ и
        $\Gamma,x:S,\Delta\vdash t:T$, то
        $\Gamma,\Delta\vdash t[x\coloneqq s]:T$.
        \item Если $\Gamma\vdash v:\sigma$ и
        $\Gamma,x:\sigma,\Delta\vdash t:T$, то
        $\Gamma,\Delta\vdash t\llbracket
        x\coloneqq v\rrbracket:T$.
        \item Если $\Gamma\vdash t:T$ и $t\leadsto t'$,
        то $\Gamma\vdash t':T$.
    \end{enumerate}
\end{theorem}
\begin{proof}
    Структурная индукция по термам и по $(\leadsto)$.
\end{proof}

\begin{theorem}[Сильная нормализуемость в~HM-R]
    Замкнутые термы, типизируемые в~HM-R, сильно
    нормализуемы относительно $(\leadsto)$ и
    редуцируются до значения.
\end{theorem}
\begin{proof}
    Применяем метод логических отношений Тейта. Пусть
    $SN$~--- множество сильно нормализуемых термов,
    редуцируемых до значения. Рассмотрим отношение
    $\mathcal{R}\subseteq
        \mathsf{PolyType}\times\mathsf{Term}$:
    \begin{itemize}
        \item $t\in\mathcal{R}_\alpha,\mathcal{R}_\top$
            тогда и только когда $t\in SN$;
        \item $t\in\mathcal{R}_{T_1\times T_2}$ тогда и
            только когда $\pi_i\,t\in\mathcal{R}_{T_i}$
            для $i=1,2$;
        \item $t\in\mathcal{R}_{S\to T}$ тогда и только
            когда для всех $s\in\mathcal{R}_S$ имеет
            место $t\,s\in\mathcal{R}_T$;
        \item $t\in\mathcal{R}_{\forall\vec{\alpha}.T}$
            тогда и только когда $t$ замкнут и
            $t\in\mathcal{R}_{
                T[\vec{\alpha}\coloneqq\vec{S}]}$
            для всех допустимых $\vec{S}$.
    \end{itemize}
    Заметим, что $\mathcal{R}_T\subseteq SN$. Более
    того, справедливо следующее свойство: если
    $t\in\mathcal{R}_{\forall\vec{\alpha}.T}$, то
    $\llbracket\mathrm{nf}(t)\rrbracket^\mathcal{M}/
    \varepsilon\in\mathcal{R}_{\llbracket T\rrbracket\,
        \mathcal{M}\,[\vec{\alpha}\coloneqq\vec{S}]}$
    для всех $\mathcal{M}$ таких, что все термы внутри
    него принадлежат соответствующим
    $\mathcal{R}_\forall$.

    Тогда мы можем доказать нашу основную лемму: если
    $\Gamma\vdash t:T$ и $\gamma$~--- подстановка
    такая, что $\gamma\,x\in\mathcal{R}_{\Gamma(x)}$
    для всех $x\in\Gamma$, то
    $\gamma\,t\in\mathcal{R}_T$.
\end{proof}

\begin{theorem}[Консервативность над~STLC]
    Нормальные формы замкнутых термов, типизируемых
    в~HM-R, типизируемы в~STLC.
\end{theorem}
\begin{proof}
    Благодаря сохранению типов нормальная форма
    $\mathrm{nf}(t)$ типизируема в~HM-R. По
    определению, она является значением; следовательно,
    не содержит подтермов переинтерпретации и
    let-связываний, а значит, типизируема в~STLC.
\end{proof}

\begin{theorem}
    Глобальный вывод типов в~HM-R разрешим.
\end{theorem}
\begin{proof}
    Достаточно расширить любой алгоритм вывода типов на
    основе унификации для системы типов
    Хиндли~---~Милнера так, чтобы он поддерживал модели
    и переинтерпретацию.
    \begin{itemize}
        \item Поскольку модели задаются явным образом,
            достаточно лишь проверить их.
        \item Поскольку полиморфные функции не зависят
            от локального контекста, их тип полностью
            определяется в момент добавления в
            контекст.
        \item После проверки модели $M$ в
            $x^\mathcal{M}$ можно напрямую вычислить
            $\llbracket T\rrbracket^\mathcal{M}_{[
            \vec{\alpha}\coloneqq\vec{\eta}]}$ (так как
            $\mathcal{M}$ известно, а $T$ полностью
            определена), где $\vec{\eta}$~--- свежие
            метaпеременные.
    \end{itemize}
\end{proof}

Таким образом, в этом разделе мы ввели корректную
систему типов на основе системы типов
Хиндли~---~Милнера, способную интерпретировать
собственные замкнутые термы в произвольной ``декартово
замкнутой категории'', определяемой внутри неё. Эта
система не является по-настоящему
\textit{рефлексивной}, поскольку модели не являются
полностью внутренними: они требуют специальной формы
операторов типов, которая недоступна остальной части
системы типов. Чтобы интернализировать определение
моделей, мы переходим в следующем разделе к системе с
зависимыми типами.

\section
[Почти рефлексивный язык с зависимыми типами: MLTT-R]
{Почти рефлексивный язык с зависимыми типами: \\
MLTT-R}
\label{MLTT-R}

Теперь мы готовы взяться за главную задачу:
спроектировать минимальный рефлексивный язык с
зависимыми типами, а также предоставить в нём
внутренние модели самого языка для полиморфизма по
вселенным, интернализации внешней параметричности и
возможности заниматься синтетической математикой внутри
него. В качестве отправной точки мы берём теорию типов
Мартина~---~Лёфа с $\Pi$-, $\Sigma$- и
$\mathrm{Id}$-типами, типом натуральных чисел и
бесконечной некумулятивной иерархией универсов. Как
обсуждалось в разделе \ref{no-go-section}, наилучшее
практическое решение проблемы переинтерпретации
здесь~--- запретить редукции под ней,
переинтерпретировать термы как есть и ограничить модели
с помощью $\mathrm{Id}$-типов. Тем не менее, по
сравнению с случаем без зависимых типов, возникает
бесконечно больше проблем, если мы хотим определить
по-настоящему рефлексивную систему:

\begin{enumerate}
    \item Если мы наивно перенесём подход к типизации
        моделей из раздела \ref{HM-R}, то для
        спецификации модели нам потребуется задать
        счётно много моделей, по одной на каждый
        универсум. И даже если мы это сделаем, к какому
        универсуму принадлежит тип моделей?
    \item Теперь, когда у нас есть вычисления в типах,
        мы обязаны разобраться с правилом конверсии под
        переинтерпретацией:
        \[
            \dfrac{t:T\quad T\equiv T'}{t:T'}
            \stackrel{?}{\leadsto}
            \dfrac{\llbracket t\rrbracket:\llbracket T
                \rrbracket\quad
                \llbracket T\rrbracket\equiv
                \llbracket T'\rrbracket
            }{\llbracket t\rrbracket:
                \llbracket T'\rrbracket}
        \]
        Проблема в том, что после интерпретации
        $\llbracket T\rrbracket$ лишь
        \textit{пропозиционально} равно
        $\llbracket T'\rrbracket$.
    \item Если мы хотим по-настоящему рефлексивную
        систему, нам нужно решить, что делать с
        \textit{переинтерпретацией переинтерпретации}:
        \[
            \lambda x.\llbracket\llbracket
                x\rrbracket_1\rrbracket_2
            \quad\text{и}\quad
            \llbracket\lambda x.\llbracket x
                \rrbracket_1\rrbracket_2
        \]
        В первом случае мы выполняем две
        последовательные интерпретации над ещё не
        известным $x$; во втором~--- интерпретируем сам
        процесс интерпретации. Что это должно значить?
\end{enumerate}

Чтобы справиться с проблемами размера универсумов,
обратим внимание на следующее наблюдение: в системе без
операторов связываний переменных по уровням каждый терм
использует лишь конечное число уровней. Следовательно,
чтобы переинтерпретировать такой терм, достаточно
задать конечное число моделей универсов, которые
совместно помещаются в некоторый более крупный универс
конечного уровня. Полиморфность универсов не нужна!
Однако теперь в суждениях о типизации нам приходится
отслеживать максимальный уровень, встречающийся в
терме.

Для интерпретации правила конверсии мы сначала
определяем, в каких именно местах оно действительно
используется; для этого \textit{необходимо}
использовать двустороннюю проверку типов. В момент
конверсии интерпретация обладает достаточной
информацией, чтобы знать как $T$, так и $T'$; таким
образом, мы можем перевести цепочку редукций,
выполняемых при проверке дефинициального равенства, в
цепочку пропозициональных равенств, взятых из условий
согласованности, доказанных в определении модели. Затем
точка конверсии обёртывается в элиминацию равенства с
синтезированным свидетелем в качестве экзаменуемого
терма.

При переинтерпретации переинтерпретации на самом деле
мы имеем дело лишь с ``двойной переинтерпретацией'':
случаи вида $\lambda x.\llbracket\llbracket x
\rrbracket_1 \rrbracket_2$ можно обработать путём
\textit{слияния} двух моделей в новую,~--- техника
напоминает \textit{Collapsing towers of interpreters}
\cite{Amin2017}; в нашей системе это не требуется,
поскольку самая внутренняя переинтерпретация в этом
случае может выполнять вычисления. Что касается
интерпретации самой интерпретации, как в
$\llbracket\lambda x.\llbracket x \rrbracket_1
\rrbracket_2$, мы явно запрещаем такие конструкции,
вводя семейство модальностей $\Box_\kappa$, означающее
``при построении этого терма не использовалась
интерпретация, а все встречающиеся в нём уровни не выше
$\kappa$''. Для переинтерпретации термы должны иметь
тип с модальностью. В нотации модальности и общей
дисциплине типизации мы в значительной мере следуем
работе Gratzer и соавторов \cite{Gratzer2019} с
некоторыми важными отличиями, которые обсуждаются ниже.
Мы также используем $\Box$ для ослабления запрета на
вычисления внутри переинтерпретации, разрешая
``вшивать'' выражения с модальностью в другие выражения
с модальностью; внутри вставок вычислениям разрешено
приводить термы к модальному виду для дальнейшей
интерпретации.

\begin{figure}
    \centering
    \begin{bnf}
        $\Gamma$ : Context ::=
        | $\cdot$ : пустой контекст
        | $\Gamma,x\,\colon T$ : расширение переменной
        | $\Gamma.\,{}_\ell\text{\faLock}$
        : фиксация контекста
        ;;
        $T$, $S$ : Type ::=
        | $\mathcal{U}_\ell$ : универсум
        | $\mathrm{El}_\ell(t)$
        : универсальное раскодирование
        | $\mathrm{Nat}_\ell$ : тип натуральных чисел
        | $(x\,\colon S)\times T$ : $\Sigma$-тип
        | $(x\,\colon S)\to T$ : $\Pi$-тип
        | $t_1=_T t_2$ : тип тождества
        | $\Box_\kappa\,T$
        : тип переинтерпретируемых термов
        | $\mathcal{M}_{\ell_0,\ldots,\ell_\kappa}$
        : тип моделей
        | $T^t_{\ell_0,\ldots,\ell_\kappa}$
        : типовая переинтерпретация
        ;;
        $\ell$, $\kappa$, $j$ : Level :in: $\mathbb{N}$
        : конечные уровни ;;
        $t$, $s$ : Term ::=
        | $x$ // $y$ // $z$ // \texttt{m} // $\ldots$
        : переменные
        | $\mathtt{U}_\ell$ : код универсума
        | $\mathtt{Nat}_\ell$
        : код типа натуральных чисел
        | $0$ : ноль
        | $\mathrm{S}\,t$ : преемник
        | $\mathrm{natind}_{x.T}\,t\,\{s_0,s_S\}$
        : индукция
        | $\mathtt{Sigma}(x\,\colon s)t$
        : код $\Sigma$-типа
        | $(s, t)$ : пара
        | $\pi_1\,t$ // $\pi_2\,t$ : проекции
        | $\mathtt{Pi}(x\,\colon s)t$
        : код $\Pi$-типа
        | $\lambda x.\,t$ : абстракция
        | $t\,s$ : применение
        | $\mathtt{Id}\,t\,t_1\,t_2$
        : код типа тождества
        | $\mathrm{refl}$ : рефлексивность
        | $\mathrm{J}_{xy.S}\,t\,s$
        : правило J (элиминация типа равенства)
        | $\mathtt{Box}_\kappa\,t$
        : код модального типа
        | $[t]_\text{\faLock}$ : фиксация
        | $[t]_\text{\faUnlock}$ : снятие фиксации
        | $\llbracket t
        \rrbracket^{\ell_0,\ldots,\ell_\kappa}$
        : внутренняя интерпретация
        | $t\,\colon T$ : аннотация типа
        ;;
    \end{bnf}
    \caption{Грамматика MLTT-R}
    \label{mltt-r-grammar}
\end{figure}

Грамматические правила, приведённые на
рис.~\ref{mltt-r-grammar}, отражают обсуждение выше.
Обратим внимание, что тип моделей $\mathcal{M}$
аннотирован последовательностью уровней
$\ell_0,\ldots,\ell_\kappa$: уровень $\ell_j$~--- это
уровень, на котором оказываются типы из $\mathcal{U}_j$
после интерпретации в такой модели.

Правила типизации также являются прямой адаптацией
общего подхода к двусторонней типизации систем с
зависимыми типами \cite{Thiago2025} и приведены на
рис.~\ref{mltt-r-wf}--\ref{mltt-r-typing}; правила
размещения по уровням построены так, что
$\Gamma\vdash_\ell\mathrm{level}(T)=\kappa$
всегда влечёт $\kappa\leqslant\ell$.

Важное отличие от стандартного подхода двусторонней
типизации, согласно которому ``конструкторы
проверяются, а элиминаторы выводятся'', состоит в том,
что уровни типов всегда выводятся: у нас нет
кумулятивности универсов и нет оператора подъёма,
поэтому мы не можем проверять уровни $\Sigma$ и $\Pi$.
Также отметим, что мы явно задаём уровни, необходимые
для проверки $t$ в типовой переинтерпретации $T^t$: это
происходит потому, что $\mathcal{M}$ дефинициально
равно некоторому огромному вложенному $\Sigma$-типу, и
сопоставление с ним было бы нецелесообразно, если бы он
выводился из $t$.

Как уже было сказано, переинтерпретация ограничена
действием на модальные типы; кроме того, мы требуем
незафиксированный контекст для переинтерпретации. В
совокупности эти ограничения гарантируют, что модальные
термы действительно не содержат подтермов
переинтерпретации. Важное отличие от
$\mathrm{MLTT}\text{\faLock}$ \cite{Gratzer2019},
помимо добавления семейства модальностей, состоит в
способе фиксации контекста в $[\_]_\text{\faLock}$ и
снятия фиксации в $[\_]_\text{\faUnlock}$: при фиксации
мы требуем, чтобы внешний контекст был незафиксирован;
при снятии фиксации, вместо удаления всех фиксаций, мы
удаляем контекст до первой фиксации. Фактически, это
накладывает ограничение, что контекст типизации никогда
не содержит более одной фиксации. Таким образом,
например, наш $\Box_\ell$ не является комонадой, но
аксиома K модальной логики всё ещё
выполняется\footnote{Интересно, что аналогичная
K-модальность впервые описана Birkedal и соавторами в
работе ``Modal dependent type theory and dependent
right adjoints'' \cite{Birkedal2020}. Совпадение: они
тоже уточняют суждения о типизации уровнями универсов!
Однако у них лишь одна модальность.}:
\[
    \dfrac{\cdot\vdash_\ell\mathrm{level}(A,B)=\_}{
        \cdot\vdash_\kappa\lambda f x.[[f
            ]_\text{\faUnlock}[x]_\text{\faUnlock}
        ]_\text{\faLock}:\Box_\ell(A\to B)\to
            \Box_\ell A\to\Box_\ell B
    }(\ell\leqslant\kappa)
\]

Таким образом, снятие фиксации по-настоящему
приобретает семантику ``вставки'' в том смысле, что оно
может содержать лишь модальные значения из ``внешнего
мира''.

\begin{figure}
    \centering
    \begin{mathpar}
        \inferrule[WF-Nil]{}{\cdot\vdash_\ell}
        \and
        \inferrule[WF-Var]{
            \Gamma\vdash_\ell\\
            \Gamma\vdash_\ell\mathrm{level}(T)=\kappa
        }{\Gamma,x:T\vdash_\ell}
        \and
        \inferrule[WF-Lock]{
            \Gamma\vdash_\ell\\
            \text{\faLock}\not\in\Gamma
        }{\Gamma.\,{}_\ell\text{\faLock}\vdash_\kappa}
        \and
        \inferrule[Level-El]
        {\Gamma\vdash_\ell t\Leftarrow
            \mathcal{U}_\kappa}
        {\Gamma\vdash_\ell\mathrm{level}(
            \mathrm{El}_\kappa(t)) = \kappa}
        \and
        \inferrule[Level-Uni]{\Gamma\vdash_\ell}{
            \Gamma\vdash_\ell
            \mathrm{level}(\mathcal{U}_\kappa) =
            \kappa+1
        }(\kappa+1\leqslant\ell)
        \and
        \inferrule[Level-Nat]{\Gamma\vdash_\ell}{
            \Gamma\vdash_\ell\mathrm{level}(
            \mathrm{Nat}_\kappa) = \kappa
        }(\kappa\leqslant\ell)
        \and
        \inferrule[Level-Sigma, Level-Pi]{
            \Gamma\vdash_\ell\mathrm{level}(S) =
            \kappa \\
            \Gamma,x:S\vdash_\ell\mathrm{level}(T) =
            j
        }{
            \Gamma\vdash_\ell
            \mathrm{level}((x:S)\times T)=
            \mathrm{level}((x:S)\to T)=
            \max(\kappa,j)
        }
        \and
        \inferrule[Level-Id]{
            \Gamma\vdash_\ell\mathrm{level}(T)=\kappa
            \\
            \Gamma\vdash_\ell t_1,t_2\Leftarrow T
        }{\Gamma\vdash_\ell\mathrm{level}(t_1=_T t_2)=
            \kappa}
        \and
        \inferrule[Level-Box]{
            \text{\faLock}\not\in\Gamma\\
            \Gamma.\,{}_\ell\text{\faLock}\vdash_\kappa
            \mathrm{level}(T)=j
        }{\Gamma\vdash_\ell\mathrm{level}(
                \Box_\kappa\,T)
            =\kappa}(\kappa\leqslant\ell)
        \and
        \inferrule[Level-Model]{\Gamma\vdash_\ell}{
            \Gamma\vdash_\ell\mathrm{level}\left(
            \mathcal{M}_{\ell_0,\ldots,\ell_\kappa}
            \right)=j
        }(j=\max(\ell_0,\ldots,\ell_\kappa
                )+2\leqslant\ell)
        \and
        \inferrule[Level-TyInterp]{
            \Gamma\vdash_\ell t\Leftarrow\mathcal{M}_{
                \ell_0,\ldots,\ell_\kappa}\\
            \text{\faLock}\not\in\Gamma\\
            \Gamma.{}_\ell\text{\faLock}\vdash_\kappa
            \mathrm{level}(T)=j
        }{\Gamma\vdash_\ell\mathrm{level}\left(
            T^t_{\ell_0,\ldots,\ell_\kappa}\right) =
            \ell_j}
    \end{mathpar}
    \caption{
        Правила корректности контекста и
        размещения по уровням}
    \label{mltt-r-wf}
\end{figure}

\begin{figure}
    \centering
    \begin{mathpar}
        \inferrule[Infer-Uni]{\Gamma\vdash_\ell}{
            \Gamma\vdash_\ell\mathtt{U}_\kappa
            \Rightarrow
            \mathcal{U}_{\kappa+1}
        }(\kappa+1\leqslant\ell)
        \and
        \inferrule[Infer-Nat]{\Gamma\vdash_\ell}{
            \Gamma\vdash_\ell\mathtt{Nat}_\kappa
            \Rightarrow
            \mathcal{U}_\kappa}(\kappa\leqslant\ell)
        \and
        \inferrule[Check-Zero]{\Gamma\vdash_\ell}{
            \Gamma\vdash_\ell 0\Leftarrow
            \mathrm{Nat}_\kappa
        }
        \and
        \inferrule[Check-Succ]
        {\Gamma\vdash_\ell t\Leftarrow
            \mathrm{Nat}_\kappa}
        {\Gamma\vdash_\ell\mathrm{S}\,t\Leftarrow
            \mathrm{Nat}_\kappa}
        \and
        \inferrule[Infer-NatInd]{
            \Gamma\vdash_\ell t\Rightarrow
            \mathrm{Nat}_\kappa\\
            \Gamma,x:\mathrm{Nat}_\kappa\vdash_\ell
            \mathrm{level}(T)=\_\\
            \Gamma\vdash_\ell s_0\Leftarrow
            T[x\coloneqq 0:\mathrm{Nat}_\kappa]\\
            \Gamma\vdash_\ell s_S\Leftarrow
            (y:\mathrm{Nat}_\kappa)\to
            (z:T[x\coloneqq y])\to
            T[x\coloneqq\mathrm{S}\,y:
            \mathrm{Nat}_\kappa]
        }{\Gamma\vdash_\ell
            \mathrm{natind}_{x.T}\,t\,\{s_0,s_S\}
            \Rightarrow T[x\coloneqq t]}
        \and
        \inferrule[Infer-Sigma]{
            \Gamma\vdash_\ell s\Rightarrow
            \mathcal{U}_\kappa\\
            \Gamma,x:\mathrm{El}_\kappa(s)\vdash_\ell t
            \Rightarrow\mathcal{U}_j
        }{\Gamma\vdash_\ell\mathtt{Sigma}(x:s)t
            \Rightarrow\mathcal{U}_{\max(\kappa,j)}}
        \and
        \inferrule[Check-Pair]{
            \Gamma\vdash_\ell s\Leftarrow S\\
            \Gamma\vdash_\ell t\Leftarrow
                T[x\coloneqq s:S]
        }{\Gamma\vdash_\ell(s,t)\Leftarrow
            (x:S)\times T}
        \and
        \inferrule[Infer-$\mathrm{Proj}_1$]
        {\Gamma\vdash_\ell t\Rightarrow(x:S)\times T}
        {\Gamma\vdash_\ell \pi_1\,t\Rightarrow S}
        \and
        \inferrule[Infer-$\mathrm{Proj}_2$]
        {\Gamma\vdash_\ell t\Rightarrow(x:S)\times T}
        {\Gamma\vdash_\ell \pi_2\,t\Rightarrow
            T[x\coloneqq \pi_1\,t]}
        \and
        \inferrule[Infer-Pi]{
            \Gamma\vdash_\ell s\Rightarrow
            \mathcal{U}_\kappa\\
            \Gamma,x:\mathrm{El}_\kappa(s)\vdash_\ell t
            \Rightarrow\mathcal{U}_j
        }{\Gamma\vdash_\ell\mathtt{Pi}(x:s)t\Rightarrow
            \mathcal{U}_{\max(\kappa,j)}}
        \and
        \inferrule[Check-Abs]
        {\Gamma,x:S\vdash_\ell t\Leftarrow T}
        {\Gamma\vdash_\ell\lambda x.\,t\Leftarrow
            (x:S)\to T}
        \and
        \inferrule[Infer-App]{
            \Gamma\vdash_\ell t\Rightarrow (x:S)\to T\\
            \Gamma\vdash_\ell s\Leftarrow S
        }{\Gamma\vdash_\ell t\,s\Rightarrow
            T[x\coloneqq s:S]}
        \and
        \inferrule[Infer-Id]{
            \Gamma\vdash_\ell t\Rightarrow
            \mathcal{U}_\kappa\\
            \Gamma\vdash_\ell t_1,t_2\Leftarrow
            \mathrm{El}_\kappa(t)
        }{\Gamma\vdash_\ell\mathtt{Id}\,t\,t_1\,t_2
            \Rightarrow\mathcal{U}_\kappa}
        \and
        \inferrule[Check-Refl]
        {\Gamma\vdash_\ell t_1\equiv t_2:T}
        {\Gamma\vdash_\ell\mathrm{refl}\Leftarrow
            t_1=_T t_2}
        \and
        \inferrule[Infer-J]{
            \Gamma\vdash_\ell t\Rightarrow t_1=_T t_2\\
            \Gamma,x:T,y:t_1=_T x\vdash_\ell
            \mathrm{level}(S)=\_\\
            \Gamma\vdash_\ell s\Leftarrow
            S[x\coloneqq t_1:T,
            y\coloneqq\mathrm{refl}:t_1=_T t_1]
        }{\Gamma\vdash_\ell\mathrm{J}_{xy.S}\,t\,s
            \Rightarrow S[x\coloneqq t_2:T,
                y\coloneqq t]}
        \and
        \inferrule[Infer-Var]{
            \text{\faLock}\not\in\Delta\\
            \Gamma,x:T,\Delta\vdash_\ell
        }{\Gamma,x:T,\Delta\vdash_\ell x\Rightarrow T}
        \and
        \inferrule[Infer-Box]{
            \text{\faLock}\not\in\Gamma\\
            \Gamma.\,{}_\ell\text{\faLock}
                \vdash_\kappa t
            \Rightarrow\mathcal{U}_j
        }{\Gamma\vdash_\ell\mathtt{Box}_\kappa\,t
            \Rightarrow\mathcal{U}_j}
        \and
        \inferrule[Check-Lock]{
            \text{\faLock}\not\in\Gamma\\
            \Gamma.\,{}_\ell\text{\faLock}
                \vdash_\kappa t
            \Leftarrow T
        }{\Gamma\vdash_\ell[t]_\text{\faLock}\Leftarrow
            \Box_\kappa T}
        \and
        \inferrule[Infer-Unlock]{
            \text{\faLock}\not\in\Delta\\
            \Gamma\vdash_\kappa t\Rightarrow\Box_j T
        }{\Gamma.{}_\kappa\text{\faLock},\Delta
            \vdash_\ell[t]_\text{\faUnlock}
            \Rightarrow T}
        \and
        \inferrule[Infer-Interp]{
            \Gamma\vdash_\ell t\Rightarrow
                \Box_\kappa T\\
            \text{\faLock}\not\in\Gamma
        }{\Gamma\vdash_\ell \llbracket t\rrbracket^{
                \ell_0,\ldots,\ell_\kappa}\Rightarrow
            \left(\mathtt{m}:\mathcal{M}_{
                \ell_0,\ldots,\ell_\kappa}\right)\to
            T^\mathtt{m}
        }
        \and
        \inferrule[Infer-Ascr]{
            \Gamma\vdash_\ell\mathrm{level}(T)=\_\\
            \Gamma\vdash_\ell t\Leftarrow T
        }{\Gamma\vdash_\ell t:T\Rightarrow T}
        \and
        \inferrule[Check-Inferred]{
            \Gamma\vdash_\ell t\Rightarrow S\\
            \Gamma\vdash_\ell S\equiv T
        }{\Gamma\vdash_\ell t\Leftarrow T}
    \end{mathpar}
    \caption{Правила двусторонней типизации}
    \label{mltt-r-typing}
\end{figure}

\subsection
{Правила вычисления MLTT-R и
модели систем с зависимыми типами}

Заметим, что приведённое выше обсуждение не зависело от
выбора формализма, в котором излагаются модели. Теперь,
однако, пришло время зафиксировать формализм, поскольку
без него невозможно дать определение того, куда
вычисляется $\mathcal{M}$.

В соответствии с современным уровнем развития области
мы используем \emph{категории с семействами}
\cite{castellan2020cwf} (или, сокращённо, CwF),
поскольку это в настоящее время один из наиболее
удобных формализмов для спецификации моделей теорий
зависимых типов. Ниже приведено самодостаточное
определение категории с семействами, которое даётся в
явном виде для удобства при последующей интернализации:

\begin{definition}[Категория]
    Категория представляет собой следующую
    структуру\footnote{Разумеется, декартово замкнутая
    категория~(\ref{ccc}) тоже является категорией.
    Часть определений повторяется для
    самодостаточности.}:
    \begin{enumerate}
        \item Множество \emph{объектов} $\mathrm{Ob}$ и
            семейство \emph{морфизмов}
            $\mathrm{Hom}(S, T)$, индексированное
            парами $S,T\in\mathrm{Ob}$, причём
            $f\in\mathrm{Hom}(S,T)$ обозначается как
            $f:S\to T$. Кроме того, объект $S$
            называется \emph{источником} морфизма $f$,
            а $T$ --- \emph{целью} морфизма $f$.
        \item Для каждого $A\in\mathrm{Ob}$ существует
            \emph{тождественный морфизм}
            $\mathrm{id}_A: A\to A$; для каждого
            $f:A\to B$ и $g:B\to C$ существует их
            \emph{композиция} $g\circ f:A\to C$.
            Если источник и цель тождественного
            морфизма ясны из контекста, подстрочный
            индекс опускается. Кроме того, выполняются
            следующие свойства:
            \begin{align*}
                \mathrm{id}\circ f = f =
                    f\circ\mathrm{id} & &
                (f \circ g) \circ h =
                    f \circ (g \circ h)
            \end{align*}
            для произвольных морфизмов $f,g,h$
            с совместимыми источниками и целями.
    \end{enumerate}
\end{definition}

\begin{definition}[Терминальный объект]
    \emph{Терминальным объектом} категории
    $\mathcal{C}$ называется объект
    $\top\in\mathrm{Ob}_\mathcal{C}$ такой, что для
    каждого $A\in\mathrm{Ob}_\mathcal{C}$ существует
    \textbf{единственный} морфизм $!_A:A\to\top$.
\end{definition}

\begin{definition}[Предпучок]
    \emph{Предпучком} над категорией $\mathcal{C}$
    (обозначается
    $P:\mathcal{C}^\mathrm{op}\to\mathrm{Set}$)
    называется отображение, которое объектам категории
    $\mathcal{C}$ ставит в соответствие множества, а
    морфизмам $\mathcal{C}$ --- функции между
    множествами, такое, что если
    $f:A\to_\mathcal{C} B$, то $P(f)$ является функцией
    из $P(B)$ в $P(A)$. При этом выполняются следующие
    свойства:
    \begin{align*}
        P(\mathrm{id}_A)=\mathrm{id}_{P(A)} & &
        P(g\circ f)=P(f)\circ P(g)
    \end{align*}
    где $\mathrm{id}_{P(A)}$ --- тождественная функция
    на множестве $P(A)$, а $P(f)\circ P(g)$ ---
    композиция функций $P(f)$ и $P(g)$.
\end{definition}

\begin{definition}[Категория элементов]
    Пусть $P:\mathcal{C}^\mathrm{op}\to\mathrm{Set}$
    --- предпучок. \emph{Категорией элементов}
    $\int_\mathcal{C}P$ называется категория, объекты
    которой являются парами $(A, x)$, где
    $A\in\mathrm{Ob}_\mathcal{C}$ и $x\in P(A)$, а
    морфизмы из $(A, x)$ в $(B, y)$ --- это морфизмы
    $f:A\to_\mathcal{C} B$ такие, что $P(f)(y)=x$.
\end{definition}

\begin{definition}[Категория с семействами]
    \label{cwfdef}
    Категория с семействами — это следующая структура:
    \begin{enumerate}
        \item \emph{Категория контекстов}
            $\mathcal{C}$, являющаяся категорией с
            терминальным объектом;
        \item \emph{Предпучок типов} $\mathrm{Ty}:
            \mathcal{C}^\mathrm{op}\to\mathrm{Set}$.
            Для данного морфизма
            $\gamma: \Gamma\to_\mathcal{C}\Delta$ и
            \emph{типа} $T\in\mathrm{Ty}(\Delta)$
            элемент $\mathrm{Ty}(\gamma)(T)\in
            \mathrm{Ty}(\Gamma)$ обычно записывают в
            сокращённой форме $T[\gamma]$;
        \item \emph{Предпучок термов} $\mathrm{Tm}:
            \left(\int_\mathcal{C}\mathrm{Ty}\right
            )^\mathrm{op}\to\mathrm{Set}$. Аналогично
            случаю с типами,
            $\mathrm{Tm}(\gamma)(t)\in\mathrm{Tm}
            (\Gamma, T)$ обычно записывают в
            сокращённой форме $t[\gamma]$;
        \item Для каждого
            $\Gamma\in\mathrm{Ob}_\mathcal{C}$ и
            $T\in\mathrm{Ty}(\Gamma)$ существует
            \emph{расширенный контекст}
            $\Gamma.T\in\mathrm{Ob}_\mathcal{C}$ вместе
            с \emph{морфизмом ослабления}
            $p:\Gamma.T\to\Gamma$ и
            \emph{элементом переменной}
            $q:\mathrm{Tm}(\Gamma.T, T[p])$. Кроме
            того, для любого $\Delta$,
            $\gamma:\Delta\to\Gamma$ и
            $t\in\mathrm{Tm}(\Delta, T[\gamma])$
            существует единственная \emph{расширенная
            подстановка}
            $\langle\gamma,t\rangle:\Delta\to\Gamma.T$,
            удовлетворяющая соотношениям:
            \begin{align*}
                p\circ\langle\gamma,t\rangle=\gamma & &
                q[\langle \gamma, t \rangle] = t
            \end{align*}
    \end{enumerate}
\end{definition}

Определение CwF (опр.~\ref{cwfdef}) в явном виде
выражает базовую структуру теории зависимых типов:
объекты категории $\mathcal{C}$ представляют контексты,
а морфизмы — подстановки, отображающие один контекст в
другой. Подстановка действует как на типы, так и на
термы; если в контексте $\Gamma$ задан корректный тип
$T$, можно расширить контекст, получив новый контекст
$\Gamma, x{:}T$, в котором существует корректно
типизированный терм $x{:}T$; подстановки в расширенный
контекст строятся путём задания подходящего для $x$
терма.

Это базовое определение легко расширяется добавлением
образователей типов (и терминов) для любой системы.
Соответственно, CwF для MLTT-R должна включать
$\mathrm{Ty}$- и $\mathrm{Tm}$-формы для аналогов
универсумов, $\Pi$- и $\Sigma$-типов, натуральных чисел
и типов тождества. Интерпретация для $\Box_\ell$ не
требуется, поскольку мы успешно предотвратили появление
типов в боксах внутри термов в боксах благодаря жёсткой
дисциплине~\faLock-ирования\footnote{Это же и есть
причина, по которой мы называем данную систему
\emph{почти} рефлексивной: её внутренняя
реинтерпретация не покрывает модальностей и саму
реинтерпретацию.}. Таким образом, контексты, требующие
интерпретации, строятся путём стандартных расширений из
пустого контекста, и обычная CwF достаточна для наших
целей. Поскольку фактическое интернализированное
определение нашей CwF в значительной мере повторяет
правила типизации и вычисления с рисунков
\ref{mltt-r-wf}--\ref{mltt-r-typing} и
\ref{mltt-r-typcomp}--\ref{mltt-r-comp}, но уже на
внутреннем языке, мы не приводим его здесь; для
получения более подробных сведений см.\ наши
дополнительные материалы.

Перейдём, наконец, к правилам вычисления, приведённым
на рисунках \ref{mltt-r-typcomp}--\ref{mltt-r-comp}
(полное отношение вычисления типов записывается как
$\Gamma\vdash_\ell T\leadsto T'$, а вычисление термов —
как $\Gamma\vdash_\ell t\leadsto t' \Leftrightarrow T$;
если тип не важен, а контекст или ограничение на
уровень просто наследуются, они опускаются ради
краткости). Благодаря интернализированному определению
CwF мы знаем, как вычисляется $\mathcal{M}$; остальные
правила вычисления типов и термов в большинстве своём
являются стандартными правилами вычисления MLTT для
двусторонней типизации (пока что без $\eta$-равенства).
Важно, что мы \textbf{не выполняем вычисление внутри
заблокированных термов и внутри типов в боксах}, чтобы
не попасть вновь под действие варианта нашей теоремы
невозможности~\ref{no-go} (модальный вариант изложен в
разделе обсуждения, см.\ раздел~\ref{global-elements}).

\begin{figure}
    \centering
    \begin{mathpar}
        \inferrule
        {t\leadsto t'\Leftarrow\mathcal{U}_\kappa}
        {\mathrm{El}_\kappa(t)\leadsto
            \mathrm{El}_\kappa(t')}
        \and
        \inferrule{}{\mathrm{El}(\mathtt{U}_\kappa)
            \leadsto\mathcal{U}_\kappa}
        \and
        \inferrule{}{\mathrm{El}(\mathtt{Nat}_\kappa)
            \leadsto\mathrm{Nat}_\kappa}
        \and
        \inferrule{}{\mathrm{El}(\mathtt{Sigma}(x:s)t)
            \leadsto(x:\mathrm{El}(s))
                \times\mathrm{El}(t)}
        \and
        \inferrule{}{\mathrm{El}(\mathtt{Pi}(x:s)t)
            \leadsto(x:\mathrm{El}(s))
                \to\mathrm{El}(t)}
        \and
        \inferrule{}{\mathrm{El}(
            \mathtt{Id}\,t\,t_1\,t_2)
            \leadsto t_1=_{\mathrm{El}(t)}t_2}
        \and
        \inferrule{}{\mathrm{El}(\mathtt{Box}_\kappa t)
            \leadsto\Box_\kappa\mathrm{El}(t)}
        \and
        \inferrule{S\leadsto S'}
        {(x:S)\times T\leadsto (x:S')\times T}
        \and
        \inferrule
        {\Gamma,x:S\vdash_\ell T\leadsto T'}
        {\Gamma\vdash_\ell (x:S)\times T\leadsto
            (x:S)\times T'}
        \and
        \inferrule{S\leadsto S'}
        {(x:S)\to T\leadsto (x:S')\to T}
        \and
        \inferrule
        {\Gamma,x:S\vdash_\ell T\leadsto T'}
        {\Gamma\vdash_\ell (x:S)\to T\leadsto
            (x:S)\to T'}
        \and
        \inferrule{T\leadsto T'}
        {t_1=_T t_2\leadsto t_1=_{T'} t_2}
        \and
        \inferrule{t_1\leadsto t_1'\Leftarrow T}
        {t_1=_T t_2\leadsto t_1'=_T t_2}
        \and
        \inferrule{t_2\leadsto t_2'\Leftarrow T}
        {t_1=_T t_2\leadsto t_1=_T t_2'}
        \and
        \inferrule{}{\mathcal{M}_{
                \ell_0,\ldots,\ell_\kappa}\leadsto
            \mathtt{modelDef}(
                \ell_0,\ldots,\ell_\kappa)}
        \and
        \inferrule{}{
            T_{\ell_0,\ldots,\ell_\kappa}^t\leadsto
            t.\mathrm{Tm}(t.\cdot, T_t^\cdot)}
    \end{mathpar}
    \caption{Вычисления на типах MLTT-R}
    \label{mltt-r-typcomp}
\end{figure}

\begin{figure}
    \centering
    \begin{mathpar}
        \inferrule{t\leadsto t'}{S\,t\leadsto S\,t'}
        \and
        \inferrule{t\leadsto t'}{
            \mathrm{natind}_{x.T}t\{s_0,s_S\}\leadsto
            \mathrm{natind}_{x.T}t'\{s_0,s_S\}
        }
        \and
        \inferrule{}{\mathrm{natind}_{x.T}
            (0:\mathrm{Nat}_\ell)\{s_0,s_S\}
            \leadsto
                (s_0:T[x\coloneqq 0:\mathrm{Nat}_\ell])
            }
        \and
        \inferrule{}{\mathrm{natind}_{x.T}
            (S\,t:\mathrm{Nat}_\ell)\{s_0,s_S\}
            \leadsto (s_S:(x:\mathrm{Nat}_\ell)\to T\to
            T[x\coloneqq S x])\,t\,(\mathrm{natind
            }_{x.T}(t:\mathrm{Nat}_\ell)\{s_0,s_S\})}
        \and
        \inferrule{s\leadsto s'}{\mathtt{Sigma}(x:s)t
            \leadsto\mathtt{Sigma}(x:s')t}
        \and
        \inferrule
        {\Gamma,x:\mathrm{El}(s)\vdash_\ell
            t\leadsto t'}
        {\Gamma\vdash_\ell\mathtt{Sigma}(x:s)t\leadsto
            \mathtt{Sigma}(x:s)t'}
        \and
        \inferrule{s\leadsto s'}{(s,t)\leadsto(s',t)}
        \and
        \inferrule{t\leadsto t'}{(s,t)\leadsto(s,t')}
        \and
        \inferrule{t\leadsto t'}
            {\pi_i t\leadsto \pi_i t'}
        \and
        \inferrule{}{\pi_i(t_1,t_2)\leadsto t_i}
        \and
        \inferrule{s\leadsto s'}{\mathtt{Pi}(x:s)t
            \leadsto\mathtt{Pi}(x:s')t}
        \and
        \inferrule
        {\Gamma,x:\mathrm{El}(s)\vdash_\ell
            t\leadsto t'}
        {\Gamma\vdash_\ell\mathtt{Pi}(x:s)t\leadsto
            \mathtt{Pi}(x:s)t'}
        \and
        \inferrule
        {\Gamma,x:S\vdash_\ell t\leadsto t':T}
        {\Gamma\vdash_\ell\lambda x.t\leadsto
            \lambda x.t'}
        \and
        \inferrule{t\leadsto t'}{t\,s\leadsto t'\,s}
        \and
        \inferrule{s\leadsto s'}{t\,s\leadsto t\,s'}
        \and
        \inferrule{\Gamma\vdash_\ell s\Rightarrow S}
        {(\lambda x.t)s\leadsto t[x\coloneqq s:S]}
        \and
        \inferrule{t\leadsto t'}
            {\mathtt{Id}\,t\,t_1\,t_2
            \leadsto\mathtt{Id}\,t'\,t_1\,t_2}
        \and
        \inferrule{t_1\leadsto t_1'}{
            \mathtt{Id}\,t\,t_1\,t_2\leadsto
            \mathtt{Id}\,t\,t_1'\,t_2
        }
        \and
        \inferrule{t_2\leadsto t_2'}{
            \mathtt{Id}\,t\,t_1\,t_2\leadsto
            \mathtt{Id}\,t\,t_1\,t_2'
        }
        \and
        \inferrule{t\leadsto t'}
            {\mathrm{J}_{xy.S}\,t\,s
            \leadsto\mathrm{J}_{xy.S}\,t'\,s}
        \and
        \inferrule{}{\mathrm{J}_{xy.S}(\mathrm{refl}:
            t_1=_T t_2)\,s\leadsto(s:S[x\coloneqq t_2,
            y\coloneqq\mathrm{refl}:t_1=_T t_2])}
        \and
        \inferrule{t\leadsto t'}{(t:T)\leadsto(t':T)}
        \and
        \inferrule{t\leadsto t'}
            {\llbracket t\rrbracket^{
                \ell_0,\ldots,\ell_\kappa}
            \leadsto\llbracket t'\rrbracket^{
                \ell_0,\ldots,\ell_\kappa}}
        \and
        \inferrule{}{\llbracket [t]_\text{\faLock}:
            \Box_\kappa T
            \rrbracket^{\ell_0,\ldots,\ell_\kappa}
                \leadsto
            \lambda\mathtt{m}.\llbracket t
                \rrbracket^\cdot_{
                \Leftarrow T}\,\mathtt{m}}
    \end{mathpar}
    \caption{Вычисления на термах MLTT-R}
    \label{mltt-r-comp}
\end{figure}

Помимо этого, необходимо также привести правила
вычисления реинтерпретации; как оказалось,
\textit{прямая реинтерпретация терминов как есть}
позволяет определить реинтерпретацию не как редукс, а
как кипу тотальных функций на типах и терминах --- так
же, как и в случае без зависимых типов
(раздел~\ref{HM-R}). Если быть точным, нам необходимы
две взаимно рекурсивные функции для термов:
$\llbracket t\rrbracket^\Gamma_{\Leftarrow T}
\mathtt{m}$ вычисляется для термов в режиме проверки, а
$\llbracket t\rrbracket^\Gamma_{\Rightarrow}\mathtt{m}$
--- для термов в режиме синтеза и возвращает выведенный
тип; обе вычисляют соответствующие элементы
$\mathrm{Tm}(\Gamma,\_)$ из CwF $\mathtt{m}$. Выражение
$T^\Gamma_\mathtt{m}$ вычисляет элемент
$\mathrm{Ty}(\Gamma)$. Для обработки правила конверсии
требуется ещё одна вспомогательная функция, вычисляющая
цепочку равенств; у неё есть две тонкости:

\begin{itemize}
    \item На данном этапе мы ещё не знаем, является ли
        система сильно нормализующейся, поэтому
        потенциально нам может встретиться бесконечная
        цепочка равенств.
    \item В системах, где равенство существенно зависит
        от доказательства (и мы хотели бы иметь
        возможность развивать теорию именно в этом
        направлении), различные цепочки редукций могут
        приводить к различным свидетелям равенства, так
        что такая функция потенциально может быть
        некорректно определённой.
\end{itemize}

Наше \emph{хитрое} решение состоит в следующем.
Зафиксируем определённую стратегию вычисления
(например, строгую) и используем её как при
двусторонней проверке типов, так и при выполнении
функции извлечения свидетелей. Если мы запускаем
функцию извлечения свидетелей на терме, который был
успешно типизирован алгоритмом типизации, то этот
алгоритм уже завершил работу, а значит, и функция
извлечения свидетелей также гарантированно завершится.
Поскольку стратегия вычисления зафиксирована, функция
извлечения свидетелей всегда будет выдавать один и тот
же результат на одном и том же входе и, следовательно,
является корректно определённой. И ещё раз: поскольку
функции интерпретации непосредственно переводят правила
двусторонней типизации термов в применения
соответствующих конструкций модели, их явное
определение не приводится.

\subsection{Метатеоретические свойства MLTT-R}
\label{MLTT-R-properties}

Ещё раз, мы ограничиваемся проверкой лишь самых
существенных свойств: конфлюэнтностью вычислений,
сохранением типов при вычислениях, завершаемостью всех
вычислений, а также тем, что получившаяся система типов
является консервативным расширением MLTT.

\begin{theorem}[Теорема Чёрча --- Россера для MLTT-R]
    Рефлексивное транзитивное замыкание отношения
    $(\leadsto)$ обладает свойством Чёрча --- Россера.
\end{theorem}
\begin{proof}
    Классическое доказательство через отношение
    параллельной редукции \cite{CRproof}. Поскольку
    интерпретация вычисляется как рекурсивная функция
    на открытых термах, она коммутирует с подстановкой,
    что является необходимым условием конфлюэнтности.
\end{proof}

\begin{theorem}[Сохранение типов для MLTT-R]
    Если $\Gamma\vdash t\Leftarrow T$ и $t\leadsto t'$,
    то $\Gamma\vdash t'\Leftarrow T$; аналогично для
    $\Gamma\vdash t\Rightarrow T$.
\end{theorem}
\begin{proof}
    Структурной индукцией по термам и $(\leadsto)$.
    Заметим, что в двусторонней системе следует
    аккуратно выстроить правила так, чтобы подстановки
    применялись к термам в режиме вывода, чтобы
    выводимые термины оставались выводимыми, а
    проверяемые --- проверяемыми.
    Аналогично доказательству сохранения типов для
    HM-R (раздел~\ref{HM-R-properties}), выполняется
    ключевая лемма:
    \[  \dfrac{
        \Gamma.{}_\ell\text{\faLock},\Delta
        \vdash_\kappa t\Leftarrow T
        \quad\Gamma\vdash_\ell s\Leftarrow
        \mathcal{M}^{\ell_0,\ldots,\ell_\kappa}
    }{\Gamma\vdash_\ell\llbracket t\rrbracket_{
            \Leftarrow T}\Rightarrow
        s.\mathrm{Tm}(s\llbracket\Delta\rrbracket,
        T_s^\Delta)}
    \]
    а также аналогичные леммы для
    $\llbracket\_ \rrbracket_\Rightarrow$ и
    $\__s^\Delta$.
\end{proof}

\begin{theorem}[Тождественная модель]
    Для каждого $\ell$ существует тождественная модель
    $I_\ell:\mathcal{M}_{0,\ldots,\ell}$ со свойством
    $\llbracket [t]_\text{\faLock}\rrbracket I_\ell
    \equiv t$ для любого замкнутого~$t$.
\end{theorem}
\begin{proof}
    В качестве типа контекстов возьмём универсум на
    соответствующем уровне; пусть
    $\mathrm{Ty}_\ell(\Gamma)=
    \Gamma\to\mathcal{U}_\ell$;
    $\mathrm{Tm}_\ell(\Gamma, T)=(\gamma:\Gamma)\to
    T\,\gamma$. Все остальные интерпретации термов
    определим так, чтобы они реинтерпретировались в те
    же самые термы.
\end{proof}

\begin{theorem}[Сильная нормализация для MLTT-R]
    \label{mltt-r-sn}
    Замкнутые термины, типизируемые в MLTT-R, сильно
    нормализуются относительно $(\leadsto)$.
\end{theorem}
\begin{proof}
    Достаточно любого доказательства нормализации для
    MLTT. Для методов, основанных на логических
    соотношениях Тейта (как было сделано для HM-R в
    разделе~\ref{HM-R-properties}), положим, что
    $t\in\mathcal{R}_{\Box_\kappa T}$ означает:
    «интерпретация $t$ принадлежит $\mathcal{R}$ для
    интерпретации $T$ для всех внутренних моделей из
    $\mathcal{R}_\mathcal{M}$». Для методов
    нормализации путём вычисления (как
    в~\cite{ABEL200717}) возьмём значение такого $t$
    равным его функции интерпретации; чтобы прочитать
    значение обратно в терм, применим эту функцию к
    тождественной модели.
\end{proof}

\begin{theorem}[Консервативность относительно MLTT]
    Для контекста~$\Gamma$ MLTT, типа~$T$ MLTT и
    терма~$t$ MLTT-R, если
    $\Gamma\vdash t\Leftarrow T$, то
    $\Gamma\vdash_\mathrm{MLTT}\mathrm{nf}(t)
        \Leftarrow T$.
\end{theorem}
\begin{proof}
    В силу сохранения типов нормальная форма
    $\mathrm{nf}(t)$ типизируется в MLTT-R. Поскольку
    $\Gamma$ не содержит ни боксов, ни блокировок,
    любой подтерм типа $\Box_\kappa T$ имеет вид
    $[\_]_\text{\faLock}$. Так как $T$ --- тип MLTT,
    каждая такая блокировка обёрнута либо в
    $[\_]_\text{\faUnlock}$, либо в
    $\llbracket\_ \rrbracket$. Во втором случае
    возникает редекс, что невозможно в нормальной
    форме; $[\_]_\text{\faUnlock}$ возможен только в
    заблокированном контексте, а значит, выше него
    находится ещё одна блокировка
    $[\_]_\text{\faLock}$. Следовательно,
    $\mathrm{nf}(t)$ не содержит подтермов типа
    $\Box_\kappa T$, а значит, не содержит и подтермов
    разблокировки и интерпретации: это терм MLTT,
    типизируемый в двустороннем MLTT (который
    получается из MLTT-R путём удаления правил,
    связанных с модальностями, блокировкой и
    реинтерпретацией).
\end{proof}

\subsection{Полезные реинтерпретации внутри MLTT-R}
\label{MLTT-R-applications}

Ниже мы приводим некоторые возможности языка, которых
можно добиться за счёт реинтерпретации без какого-либо
дополнительного расширения языка.

\subsubsection
{Примитивная, но эффективная стратификация}
\label{MLTT-R-displacement}

По Макбрайду
\cite{McBride:Crude2011,McBride:Crude2012}, техника
примитивной, но эффективной стратификации работает
следующим образом: вместо предоставления полной
поддержки полиморфизма уровней вводим \emph{оператор
смещения}, который, получив терм $t$, использующий
универсумы $U_i,U_j,U_k,\ldots$, порождает терм, равный
$t$ за исключением того, что универсумы заменены на
$U_{i+1},U_{j+1},U_{k+1},\ldots$. Смещённый терм также
корректно типизирован, однако определяется <<на один
уровень выше>>. Хотя его нельзя непосредственно
определить как единый оператор внутри MLTT-R, для
каждого $\kappa,\ell$ можно определить операторы
смещения $\uparrow_\ell^\kappa:\mathcal{M}_{\kappa,
\kappa+1,\ldots,\kappa+\ell}$, которые сдвигают на
$\kappa$ уровней вверх закоробленное определение,
использующее универсумы вплоть до $\mathcal{U}_\ell$. В
этой модели базовые предпучки категории контекстов
смещены на $\kappa$ по сравнению с тождественной
моделью; все остальные интерпретации термов совпадают
с определениями в тождественной модели.

Немного подробностей об определении
$\uparrow_\ell^\kappa$ приведено в дополнительных
материалах.

\subsubsection{Внутренняя внешняя параметричность}
\label{MLTT-R-parametricity}

Как показано Рейнольдсом \cite{reynolds1983types},
параметричность — это свойство системы типов, которое
можно выразить простыми словами так: <<полиморфная
функция не может существенным образом зависеть от своих
параметров>>. Таким образом, любая полиморфная функция
типа $\sigma$ удовлетворяет некоторой <<свободной
теореме>> $\llbracket \sigma \rrbracket$, определяемой
посредством реляционной интерпретации типа.

В недавних работах \cite{Bernardy:Parametricity2010,
Altenkirch:Parametricity2024} предложено упрощённое
определение интерпретации параметричности для систем с
зависимыми типами, названное <<внешней
параметричностью>>. И вновь: эта реляционная
интерпретация определима в нашей системе
\emph{внутренним} образом с помощью реляционной модели
$\mathrm{Para}_\ell:\mathcal{M}_{1,\ldots,\ell+1}$.
Смещение на один уровень обусловлено тем, что тип
$T:\mathcal{U}_\ell$ в такой модели интерпретируется
тройкой $(T_1\,T_2:\mathcal{U}_\ell)\times
(T_1\to T_2\to\mathcal{U}_\ell)$. Поскольку
интерпретируемые функции не зависят от локального
контекста, мы обходим сложности с реляционным
контекстом, отмеченные Альтенкирхом и соавторами
\cite{Altenkirch:Parametricity2024}; в результате
интерпретация даёт полноценное доказательство свободной
теоремы.

Немного подробностей об определении
$\mathrm{Para}_\ell$ приведено в дополнительных
материалах.

\subsubsection{Синтетическая математика}
\label{MLTT-R-synthetic}

Как подробно обсуждалось в недавней статье Навруцкого и
соавторов \cite{Nawrocki:Lean2026}, благодаря тому, что
MLTT служит внутренним языком определённых
$(\infty,1)$-категорий, в нём можно вести
\emph{синтетическую математику}: доказать внутри MLTT
теорему о типах, интерпретировать её в модели MLTT и
тем самым получить доказательство об объектах этой
модели. Легко видеть, что оператор реинтерпретации,
имеющийся в MLTT-R, позволяет вести синтетическую
математику непосредственно в ней и сразу применять
результаты, интерпретированные в подходящей модели.

\section{Обсуждение и связанные работы}

\subsection{Что было сделано до нас}

Первоначальная работа Ламбека \cite{Lambek1980} задала
тон последующим исследованиям, установив соответствие
между системами типов, основанными на
лямбда-исчислении, и различными видами категорий,
начиная с классического соответствия между просто
типизированным лямбда-исчислением и декартово
замкнутыми категориями.

В статье Конола Эллиота «Компиляция в категории»
\cite{Elliott:Compcat2017} показано, как с помощью
плагина компилятора GHC интерпретировать (часть)
Haskell в произвольной категории (опять же,
определяемой в Haskell). Эта статья является прямым
источником вдохновения для данной работы. Впрочем, есть
важное различие: реинтерпретация в наших системах типов
является \emph{внутренней}, что открывает новые способы
её использования, в частности (но не только)
полиморфизм универсумов, внешнюю параметричность и
многое другое. Вместе с тем подход Конола более гибок:
его трансляция в категории требует от целевой категории
реализации лишь тех примитивов, которые фактически
используются при трансляции данного конкретного терма.

В работе Конора Макбрайда «Примитивная, но эффективная
стратификация»
\cite{McBride:Crude2011,McBride:Crude2012}
(систематическое изложение этой идеи дано Фавонией и
соавторами \cite{Favonia:Displacement2023}) предложена
прорывная идея использования \emph{смещения}
(displacement) как лёгкой формы полиморфизма
универсумов: чтобы повторно использовать термин,
типизированный на более низком уровне универсумов,
достаточно поднять все уровни в терме, добавив
постоянное смещение $s$. С точки зрения, продвигаемой
нами, это можно рассматривать как реинтерпретацию
выражения во внутренней модели системы типов, в которой
универсум на уровне $i$ интерпретируется как универсум
на уровне $i+s$ (раздел~\ref{MLTT-R-displacement}).

В статье Жана-Филиппа Бернарди, Патрика Йанссона и
Росса Патерсона «Параметричность и зависимые типы»
\cite{Bernardy:Parametricity2010} авторы предлагают
прямую процедуру трансляции, позволяющую получать
теоремы и доказательства параметричности из полиморфных
определений в системе с зависимыми типами. Однако
определение <<внешней параметричности>>, которому мы
следуем, дано Торстеном Альтенкирхом, Йорго Шамоуном,
Амбрусом Капоши и Майклом Шульманом в статье
«Внутренняя параметричность без интервала»
\cite{Altenkirch:Parametricity2024}. И вновь мы можем
переформулировать предложенную ими трансляцию как
реинтерпретацию во внутренней модели. Альтенкирх и
соавторы указали на проблему данного подхода: если
применить внешнюю трансляцию к суждению с непустым
контекстом, контекст тоже становится параметрическим;
именно поэтому они добавляют операцию \texttt{rel} и
получают систему типов с вычислительными свойствами
\emph{высшей наблюдательной теории типов}. Мы обходим
проблему с контекстом, поскольку реинтерпретация
применяется только к термам, контекст которых тоже
<<реинтерпретируем>> и не затрагивается трансляцией
(раздел~\ref{MLTT-R-parametricity}).

\subsection{Что можно было сделать иначе}

\subsubsection
{Просто использовать метапрограммирование}

В недавней статье «Сертифицирующее средство
интерактивных доказательств для синтетической
математики в Lean» Войцех Навруцкий и соавторы
\cite{Nawrocki:Lean2026} адаптировали подход
\emph{«Компиляция в категории»}
\cite{Elliott:Compcat2017} к системе Lean с зависимыми
типами, используя категории с семействами для
представления внутренних моделей и применяя мощные
средства метапрограммирования Lean вместо плагина
компилятора. В результате им удалось создать основу
для ведения синтетической математики в Lean:
доказательства, полученные в MLTT, применяются к более
специфическим категориям. Таким образом, наши
результаты в разделе~\ref{MLTT-R-synthetic} близки к
результатам, показанным в их статье. Более того, их
реализация уже встроена в практическую реализацию
средства интерактивных доказательств. Тем не менее,
наша система типов не использует метапрограммирование,
поскольку реинтерпретация является внутренней для неё,
и она корректно обрабатывает правило конверсии. К тому
же акценты немного различаются: мы используем
реинтерпретацию для выражения фундаментальных языковых
возможностей.

\subsubsection{
    Вычислительная модальность глобальных элементов}
\label{global-elements}

Поскольку в разделе \ref{MLTT-R} мы ввели семейство
модальностей $\Box_\ell$ реинтерпретируемых терминов,
можно было бы надеяться допустить вычисление внутри
закоробленных терминов. Чуть обобщая, следующее
определение пригодится далее:

\begin{definition}
    Система типов с зависимыми типами и типом функций
    $(\to)$ обладает \textit{вычислительной
    модальностью глобальных элементов} $\Box$, если
    существуют терм с одним пропуском
    $\mathrm{nec}\,\_$ и атомарный терм с двумя
    пропусками $\mathrm{K}\,\_\,\_$ такие, что
    допустимы следующие правила:
    \begin{mathpar}
        \inferrule{\cdot\vdash T\,\mathrm{type}}
        {\Gamma\vdash\Box\,T\,\mathrm{type}}
        \and
        \inferrule{\cdot\vdash t:T}
        {\Gamma\vdash\mathrm{nec}\,t:\Box\,T}
        \and
        \inferrule{
            \Gamma\vdash f:\Box\,(S\to T)\\
            \Gamma\vdash t:\Box\,S
        }{\Gamma\vdash K\,f\,t:\Box\,T}
    \end{mathpar}
    а также допустима редукция:
    $K\,(\mathrm{nec}\,\lambda x.t)\,
    (\mathrm{nec}\,s)\leadsto^*
    \mathrm{nec}(t[x\coloneqq s])$.
\end{definition}

За исключением вычислительного правила, эти свойства в
точности совпадают со свойствами модальностей
$\Box_\ell$ из исходного раздела~\ref{MLTT-R}. Интуиция
за данной редукцией следующая: если мы применим
«глобальную функцию» к «глобальному аргументу», то мы
должны иметь возможность «встроить» это применение, ---
что в точности и делает данная редукция.

Тогда соответствующее ослабление вычислительной
реинтерпретации можно определить следующим образом.

\begin{definition}
    [Модальная вычислительная реинтерпретация]
    Оператор типов $F[\_]$ называется \emph{модальной
    вычислительной реинтерпретацией типов}, если он
    определён на закоробленных типах (если
    $\Gamma\vdash\Box\,T\,\mathrm{Type}$, то
    $\Gamma\vdash F[T]\,\mathrm{Type}$), сохраняет
    вычисления на типах (если $T\to T'$, то
    $F[T]\to^* F[T']$) и вычисляется как примитивно
    рекурсивная функция на нормальных формах типов.
    Аналогично, терм с пропуском
    $\llbracket\_\rrbracket$ называется \emph{модальной
    вычислительной реинтерпретацией термов}, если он
    сохраняет вычисления на термах, вычисляется как
    примитивно рекурсивная функция на нормальных формах
    термов и, если $\Gamma\vdash t:\Box\,T$, то
    $\Gamma\vdash\llbracket t\rrbracket:F[T]$.
\end{definition}

Возможно, такое определение реинтерпретации
«правильно»? К сожалению, для этого случая выполняется
теорема о невозможности, аналогичная наивной
версии~\ref{no-go}:

\begin{theorem}
    [Теорема о невозможности модальной реинтерпретации]
    В конфлюентной системе типов с модальной
    вычислительной реинтерпретацией для произвольных
    корректно типизированных $s$ и $t$ должно
    выполняться:
    \[
    \mathtt{app}\llbracket\mathrm{nec}\,\lambda x.t
    \rrbracket\llbracket\mathrm{nec}\,s\rrbracket
    \equiv\llbracket\mathrm{nec}(t[x\coloneqq s])
    \rrbracket
    \]
    где $\equiv$ --- дефинициальное равенство, а
    $\mathtt{app}$ --- результат действия
    $\llbracket\_\rrbracket$ на $\mathrm{K}$.
\end{theorem}

\begin{proof}
    На конкретном примере. Пусть заданы корректные типы
    $\cdot\vdash S,T\,\mathrm{type}$ и корректно
    типизированные термы $\cdot\vdash s:S$ и
    $\cdot,x:S\vdash t:T$. Рассмотрим следующий терм:
    \[
    \Gamma,x:\Box(S\to T),y:\Box\,S\vdash
    \llbracket K\,x\,y\rrbracket:F[T]
    \]
    Поскольку $\mathrm{K}$ атомарен, терм под
    интерпретацией находится в нормальной форме, а
    значит весь термин вычисляется в
    $\mathtt{app}\llbracket x\rrbracket
    \llbracket y\rrbracket$. Подставив
    $\mathrm{nec}\,\lambda x.t$ вместо $x$ и
    $\mathrm{nec}\,s$ вместо $y$, получим
    $\mathtt{app}\llbracket\mathrm{nec}\,
    \lambda x.t\rrbracket\llbracket\mathrm{nec}\,
    s\rrbracket$.

    С другой стороны, если сначала выполнить
    подстановку, то получим $\Gamma\vdash\llbracket K\,
    (\mathrm{nec}\,\lambda x.t)(\mathrm{nec}\,s)
    \rrbracket:F[T]$. Поскольку модальная
    реинтерпретация сохраняет редукцию, этот термин
    редуцируется в
    $\llbracket\mathrm{nec}(t[x\coloneqq s])
    \rrbracket$.

    В конфлюентной системе типов с функциями
    многошаговая редукция переставляется с
    подстановкой, поэтому два термина дефинициально
    эквивалентны.
\end{proof}

Разумеется, это не запрещает нам использовать такую
модальность: как и в случае наивной интерпретации,
данный подход работает, если нам удаётся ограничить
модели теми, в которых равенства выполняются
дефинициально.

Ещё одним преимуществом модальной вычислительной
реинтерпретации является следующее свойство:

\begin{theorem}
    [Схлопывание вычислительных модальностей]
    Для MLTT-R справедливо следующее утверждение:
    \[\dfrac{
        \exists\kappa:
        \Gamma.\,{}_\ell\text{\faLock}\vdash_\kappa
        t:T
    }{\forall j\geqslant\mathrm{level}(T):
        \Gamma.\,{}_\ell\text{\faLock}\vdash_j
        \mathrm{nf}(t):\mathrm{nf}(T)}\]
\end{theorem}

Прямое доказательство здесь не приводится, однако общая
идея в том, что любой уровень, превышающий уровень $T$,
может появляться лишь в каком-то редексе $t$ или $T$.
Тогда, если бы модальности были вычислительными,
достаточно было бы использовать единственную
модальность $\Box$ вместо семейства модальностей
$\Box_\ell$: уровень типа под боксом определял бы,
сколько моделей универсумов нужно предоставить.

Но, снова, всё сводится к определению типа
\emph{дефинициальных} внутренних моделей.

\subsubsection{Реинтерпретация реинтерпретации}

В разделе \ref{MLTT-R} мы ввели нашу модальность
``реинтерпретируемых выражений'' довольно жёстко ровно
затем, чтобы не приходилось давать интерпретацию этой
модальности внутри внутренних моделей. Опять-таки,
проблема возникает с выражениями, которые пытаются
интерпретировать переменные из контекста CwF;
простейший пример — $\llbracket \lambda x.\llbracket x
\rrbracket_1\rrbracket_2$. Чтобы суметь
интерпретировать такое выражение, нам придётся
определить \textit{что оператор реинтерпретации значит
внутри внутренней модели}. Это быстро выходит из-под
контроля: вместо того чтобы просто превратить тип
моделей в (рекурсивный? индуктивный?) тип, тип моделей
превратился бы в тип, в котором тип одного из его полей
зависит от значения всей модели. На деле такое уже
случалось: Коломатская и
др.~\cite{Kolomatskaia_Shulman_2025} уже определяли тип
такого рода внутри своей \textit{теории отображённых
типов}: тип SST полу-симплициальных типов определён
коиндуктивно через значения самого себя\footnote{В
более общем случае Шульман называет такие типы
\textit{модальными коданными}\cite{shulman2026}.}. Хотя
аналогичное определение типа моделей могло бы стать
интересным расширением системы типов, представленной в
этой статье, на данный момент мы считаем такой подход
выходящим за рамки нашей работы.

\section{Заключение и дальнейшая работа}

В ретроспективе, нам удалось определить две
рефлексивные системы типов, STLC-R и MLTT-R, то есть
они обладают понятием собственных ``внутренних
моделей'', используя которые можно перетолковывать
термы самой системы внутри таких моделей. Благодаря
этому такие системы типов обретают новые возможности,
которые ранее задавались как отдельные встроенные
особенности системы. Основные метатеоретические
свойства доказаны; на будущее мы запланировали
следующее:

\begin{itemize}
    \item Возможно, стоит сделать дисциплину блокировки
        в разделе~\ref{MLTT-R} более гибкой, допуская
        более одной блокировки в контексте. С высокой
        долей вероятности это потребует включить
        модальности в определение типа ``внутренней
        модели'' и приблизит нас к по-настоящему
        рефлексивной системе.
    \item Наша техника реинтерпретации правила
        конверсии в разделе~\ref{MLTT-R} работает, но
        текущий подход к извлечению свидетеля равенства
        представляет собой \textit{хак}. Корректное
        решение потребовало бы задать извлечение
        свидетеля как редуцируемое выражение и,
        возможно, ввести определённые \textit{алгебры
        стратегий редукции}, чтобы сделать этот процесс
        более управляемым в системах, где свидетелей
        равенства может быть несколько.
    \item Подход с категориями с семействами для
        задания внутренних моделей здесь, в
        определённом смысле, представляет собой
        ``глубокий'' интерфейс: хотя все его средства
        корректно используются при реинтерпретации,
        легко добавить возможность интерпретировать
        выражения в произвольном контексте CwF. Это
        указывает на естественное расширение нашей
        модальности реинтерпретируемых выражений до
        полноценных контекстуальных типов.
    \item После реализации контекстуальных типов
        становится возможным обрабатывать ``частичную''
        реинтерпретацию выражений: если выражение,
        находящееся в процессе интерпретации, не
        использует определённую возможность системы
        типов, то не должно быть необходимости
        указывать её в модели! Это приблизит наш язык к
        подходу Конала Эллиота ``Компиляция в
        категории'' \cite{Elliott:Compcat2017} и к
        \textit{настоящему} формальному языку для
        синтетической математики.
    \item Было бы хорошо сделать равенство в MLTT-R
        более экстензиональным, скорее всего,
        приблизив систему типов к гомотопической теории
        типов \cite{hottbook}.
    \item В определённый момент нам придётся проявить
        должную тщательность: предоставить
        формализованные доказательства
        метатеоретических свойств систем и реализацию
        MLTT-R в практичном инструменте верификации
        доказательств.
\end{itemize}
